\documentclass[11pt,a4paper]{article}

\usepackage[T1]{fontenc}
\usepackage[utf8]{inputenc}
\IfFileExists{libertinus.sty}{\usepackage{libertinus}}{\usepackage{newtxtext,newtxmath}}
\usepackage[final]{microtype}
\usepackage[a4paper,margin=26mm,headheight=15pt]{geometry}
\usepackage{amsmath,amssymb,amsthm,mathtools,mathrsfs,bm}
\usepackage{booktabs,longtable,tabularx,array,multirow}
\usepackage{graphicx}
\usepackage{enumitem}
\usepackage{xcolor}
\usepackage[most]{tcolorbox}
\usepackage{listings}
\usepackage{fancyhdr}
\usepackage{xurl}
\usepackage{hyperref}
\usepackage[nameinlink,capitalise,noabbrev]{cleveref}

\definecolor{linkblue}{RGB}{24,72,124}
\definecolor{deepgreen}{RGB}{23,102,69}
\definecolor{deepred}{RGB}{140,42,42}
\definecolor{deepgold}{RGB}{122,91,19}
\definecolor{shadegray}{RGB}{246,247,249}
\definecolor{rulegray}{RGB}{195,201,208}
\definecolor{palered}{RGB}{251,236,236}
\definecolor{palegold}{RGB}{251,245,230}
\definecolor{palegreen}{RGB}{233,245,243}
\definecolor{paleblue}{RGB}{237,244,252}

\hypersetup{
  hypertexnames=false,
  colorlinks=true,
  linkcolor=linkblue,
  citecolor=deepgreen,
  urlcolor=linkblue,
  pdftitle={Inverting the Subfactorial at Infinity},
  pdfauthor={Research note prepared for the ProveIt project},
  pdfsubject={Bell-sector transseries for inverse subfactorial functions and a general rapid-core reversion calculus},
  pdfkeywords={subfactorial, derangement, inverse function, transseries, Bell numbers, inverse gamma, Lambert W, Lagrange-Burmann inversion}
}

\pagestyle{fancy}
\fancyhf{}
\fancyhead[L]{\small Inverting the subfactorial at infinity}
\fancyhead[R]{\small\nouppercase{\leftmark}}
\fancyfoot[C]{\thepage}
\renewcommand{\headrulewidth}{0.4pt}

\setlength{\parindent}{0pt}
\setlength{\parskip}{0.55em}
\setcounter{tocdepth}{2}
\setcounter{secnumdepth}{3}
\setlist[itemize]{leftmargin=1.55em,itemsep=0.25ex,topsep=0.55ex}
\setlist[enumerate]{leftmargin=1.9em,itemsep=0.35ex,topsep=0.55ex}
\allowdisplaybreaks[3]
\numberwithin{equation}{section}

\newtheorem{theorem}{Theorem}[section]
\newtheorem{proposition}[theorem]{Proposition}
\newtheorem{lemma}[theorem]{Lemma}
\newtheorem{corollary}[theorem]{Corollary}
\newtheorem{conjecture}[theorem]{Conjecture}
\theoremstyle{definition}
\newtheorem{definition}[theorem]{Definition}
\newtheorem{example}[theorem]{Example}
\newtheorem{algorithm}[theorem]{Algorithm}
\theoremstyle{remark}
\newtheorem{remark}[theorem]{Remark}

\newtcolorbox{mainresultbox}[1][]{
  enhanced,breakable,colback=paleblue,colframe=linkblue,
  boxrule=0.75pt,arc=1.2mm,left=2.5mm,right=2.5mm,top=2mm,bottom=2mm,
  title={Main result},fonttitle=\bfseries,#1}
\newtcolorbox{summarybox}[1][]{
  enhanced,breakable,colback=palegold,colframe=deepgold,
  boxrule=0.6pt,arc=1.2mm,left=2.2mm,right=2.2mm,top=1.6mm,bottom=1.6mm,
  title={Structural summary},fonttitle=\bfseries,#1}
\newtcolorbox{warningbox}[1][]{
  enhanced,breakable,colback=palered,colframe=deepred,
  boxrule=0.6pt,arc=1.2mm,left=2.2mm,right=2.2mm,top=1.6mm,bottom=1.6mm,
  title={A branch issue that cannot be ignored},fonttitle=\bfseries,#1}
\newtcolorbox{checkpointbox}[1][]{
  enhanced,breakable,colback=palegreen,colframe=deepgreen,
  boxrule=0.6pt,arc=1.2mm,left=2.2mm,right=2.2mm,top=1.6mm,bottom=1.6mm,
  title={Computational checkpoint},fonttitle=\bfseries,#1}

\lstdefinestyle{pseudo}{
  basicstyle=\ttfamily\small,columns=fullflexible,frame=single,
  rulecolor=\color{rulegray},backgroundcolor=\color{shadegray},
  showstringspaces=false,keepspaces=true,breaklines=true,
  xleftmargin=0.4em,xrightmargin=0.4em,aboveskip=0.8em,belowskip=0.8em}
\lstset{style=pseudo}

\newcommand{\e}{\mathrm e}
\newcommand{\ii}{\mathrm i}
\newcommand{\dd}{\,\mathrm d}
\newcommand{\Ord}{\mathcal O}
\newcommand{\littleo}{\mathrm o}
\newcommand{\Bell}{\mathrm B}
\newcommand{\Sub}{\mathfrak D}
\newcommand{\Core}{\mathcal A}
\newcommand{\Tail}{\mathcal J}
\newcommand{\InvCore}{\mathfrak a}
\newcommand{\Lam}{\Lambda}
\newcommand{\Res}{\mathcal R}
\newcommand{\med}{\mathrm{med}}
\newcommand{\principal}{\mathrm{pr}}
\newcommand{\coeff}[2]{[\,#1^{#2}\,]}
\newcommand{\falling}[2]{#1^{\underline{#2}}}
\DeclareMathOperator{\Resi}{Res}
\DeclareMathOperator{\sgn}{sgn}
\DeclareMathOperator{\Arg}{Arg}
\DeclareMathOperator{\Li}{Li}

\title{\textbf{Inverting the Subfactorial at Infinity}\\[0.35em]
\Large Bell-Sector Transseries, Inverse-Gamma Geometry,\\
and a General Reversion Calculus for Rapid Cores with Tiny Oscillatory Tails}
\author{Research note prepared for the ProveIt project}
\date{3 September 2026}

\begin{document}
\maketitle

\begin{abstract}
The integer subfactorial $D_n={!n}$ is asymptotic to $n!/\e$, but its
exponentially small absolute correction has a complete inverse-power expansion
whose coefficients are Bell numbers.  Inverting this relation is subtler than
inverting Stirling's formula: a noninteger subfactorial requires a choice of
analytic continuation, the dominant inverse is the principal inverse gamma
function, and the Bell correction creates a new family of sectors beyond the
ordinary Lambert--Stirling scale.

We begin with the lateral analytic continuations
\[
  \Sub_{\sigma}(z)=\frac{\Gamma_{\sigma}(z+1,-1)}{\e},
  \qquad \sigma\in\{+1,-1\},
\]
where the negative argument is reached with argument $\sigma\pi$.  For
$\Re z>-1$ they admit the exact decomposition
\[
  \Sub_{\sigma}(z)=\frac{\Gamma(z+1)}{\e}
  +\e^{\ii\sigma\pi z}\Tail(z),
  \qquad
  \Tail(z)=\frac1\e\int_0^1 \e^t t^z\dd t
          =\frac1\e\sum_{j\ge0}\frac1{j!(z+j+1)}.
\]
The endpoint integral satisfies
\[
  \Tail(x)\sim\sum_{k\ge1}(-1)^{k-1}\frac{\Bell_k}{x^k},
\]
and we give an exact signed remainder whose absolute value is bounded by
$\Bell_{m+1}x^{-m-1}$ after $m$ terms.  Thus the Bell expansion is not merely
an integer identity; it is the slow coefficient expansion of an analytic
oscillatory tail.

Let $\InvCore(y)$ be the large real solution of
$\Gamma(\InvCore(y)+1)/\e=y$, and put $P(x)=\psi(x+1)$.  There is, for each
large positive $y$, a unique inverse branch $X_{\sigma}(y)$ of
$\Sub_{\sigma}$ satisfying $X_{\sigma}(y)-\InvCore(y)\to0$.  Its complete
Bell-sector transseries is
\[
 X_{\sigma}(y)\sim \InvCore(y)
 +\sum_{n\ge1}\frac{\e^{\ii\sigma n\pi\InvCore(y)}}{y^n}
 C_{n,\sigma}(\InvCore(y)),
\]
where
\[
 C_{n,\sigma}(x)=\frac{(-1)^n}{n!}
 \bigl(\mathcal L^{(\sigma)}_{n-1,n}\circ\cdots\circ
       \mathcal L^{(\sigma)}_{1,n}\bigr)
 \left(\frac{\Tail(x)^n}{P(x)}\right),
 \quad
 \mathcal L^{(\sigma)}_{r,n}f
 =\frac{f'}{P}+\left(\frac{\ii\sigma n\pi}{P}-r\right)f.
\]
We calculate the first sectors explicitly, expand their coefficients in Bell
numbers and logarithms, prove fixed-sector error estimates, and verify the
formula numerically.  The core $\InvCore$ is itself expanded to arbitrary
order by a centered Stirling inversion.  If
$\lambda=\log(\e y/\sqrt{2\pi})$, $w=W_0(\lambda/\e)$, and
$u_0=\lambda/w$, then
\[
 \InvCore(y)=-\frac12+u_0+\sum_{m\ge1}
 \frac{p_m(w)}{u_0^{2m-1}},
\]
with an explicit triangular recurrence for the rational functions $p_m$.

The final part abstracts the calculation.  For a rapidly varying core
$\Core(x)$ and a small slow or oscillatory tail $\Res(x)$, we prove the
all-orders identity
\[
 (\Core+\Res)^{-1}(y)\sim a(y)+
 \sum_{n\ge1}\frac{(-1)^n}{n!}
 \frac{\dd^{n-1}}{\dd y^{n-1}}
 \left[a'(y)\Res(a(y))^n\right],
 \qquad a=\Core^{-1},
\]
and convert it into a valuation-raising operator calculus.  This supplies a
reusable transseries inversion theory for factorial, gamma-type, Barnes-type,
and endpoint-Mellin perturbations.  We also treat a real median interpolation
and the genuinely discrete generalized inverse, whose staircase nature must be
kept separate from any smooth inverse transseries.
\end{abstract}

\begin{mainresultbox}
Let
\[
 \Core(x)=\frac{\Gamma(x+1)}\e,
 \qquad
 \Tail(x)=\frac1\e\int_0^1\e^t t^x\dd t,
 \qquad
 \Core(\InvCore(y))=y,
 \qquad P(x)=\psi(x+1).
\]
For $y\to+\infty$, the two lateral inverse branches have the sector-exact
normal form
\begin{align*}
X_{\sigma}(y)={}&\InvCore(y)
 -\frac{\e^{\ii\sigma\pi\InvCore(y)}}{y}
     \frac{\Tail(\InvCore(y))}{P(\InvCore(y))}\\
&+\frac{\e^{2\ii\sigma\pi\InvCore(y)}}{y^2}
 \left[
 \frac{\Tail\Tail'}{P^2}
 +\frac{\ii\sigma\pi\Tail^2}{P^2}
 -\frac{\Tail^2P'}{2P^3}
 -\frac{\Tail^2}{2P}
 \right]_{x=\InvCore(y)}\\
&+\Ord\!\left(
 \frac{1}{y^3\InvCore(y)^3\log\InvCore(y)}
 \right).
\end{align*}
Writing $a=\InvCore(y)$ and $\ell=\log a$, the first sector begins as
\[
 -\frac{\e^{\ii\sigma\pi a}}{y}
 \left[\begin{aligned}
 &\frac1{a\ell}
 -\frac{4\ell+1}{2a^2\ell^2}
 +\frac{60\ell^2+13\ell+3}{12a^3\ell^3}\\[-0.2em]
 &\hspace{1.2em}
 -\frac{360\ell^3+64\ell^2+14\ell+3}{24a^4\ell^4}
 +\Ord\!\left(\frac1{a^5\ell}\right)
 \end{aligned}\right]
\]
and the second begins as
\[
 \frac{\e^{2\ii\sigma\pi a}}{y^2}
 \left[
 -\frac1{2a^2\ell}+\frac{\ii\sigma\pi}{a^2\ell^2}
 +\frac{8\ell^2-3\ell-2-16\ii\sigma\pi\ell-4\ii\sigma\pi}
 {4a^3\ell^3}
 +\Ord\!\left(\frac1{a^4\ell}\right)
 \right].
\]
At general sector number $n$,
\[
 C_{n,\sigma}(a)\sim-\frac1n\frac{\Tail(a)^n}{P(a)}
 \sim-\frac1{n a^n\log a}.
\]
The core is
\begin{align*}
\InvCore(y)={}&-\frac12+u_0
 +\frac{1}{24(1+w)u_0}
 -\frac{14w^2+38w+29}{5760(1+w)^3u_0^3}\\
&+\frac{2232w^4+10104w^3+17634w^2+14234w+4577}
 {2903040(1+w)^5u_0^5}
 +\Ord(u_0^{-7}),
\end{align*}
where $\lambda=\log(\e y/\sqrt{2\pi})$, $w=W_0(\lambda/\e)$, and
$u_0=\lambda/w$.
\end{mainresultbox}

\begin{warningbox}
The integer sequence $D_n$ has no ordinary inverse function.  Moreover, the
standard incomplete-gamma continuation is complex-valued between the
integers.  Consequently, three distinct objects must not be conflated:
(i) a lateral complex inverse branch $X_{\sigma}$; (ii) a chosen real
interpolation and its inverse; and (iii) the integer-valued generalized
inverse $\min\{n:D_n\ge y\}$.  All three are analyzed below, but they have
different asymptotic meanings.
\end{warningbox}

\tableofcontents
\newpage

\section{Introduction and orientation}

For an integer $n\ge0$, the subfactorial or derangement number
\begin{equation}\label{eq:Dn-def}
 D_n={!n}=n!\sum_{j=0}^{n}\frac{(-1)^j}{j!}
\end{equation}
counts permutations of $n$ objects having no fixed point.  Its leading
asymptotic behavior is the elementary relation $D_n\sim n!/\e$.  A much more
informative formula, emphasized by Hassani~\cite{Hassani2020}, is
\begin{equation}\label{eq:Hassani-intro}
 D_n=\frac{n!}{\e}
 +\sum_{k=1}^{m}(-1)^{n+k-1}\frac{\Bell_k}{n^k}
 +\Ord(n^{-m-1}),
\end{equation}
where $m$ is fixed, $\Bell_k$ is the $k$th Bell number, and the constant in
the remainder can be taken no larger than $\Bell_{m+1}$.

The purpose of this article is to reverse \eqref{eq:Hassani-intro}.  The
naive instruction ``solve $D_x=y$'' hides four layers of structure.

\begin{enumerate}
\item The integer sequence has to be embedded in a continuous or analytic
      object before it has a local inverse in the usual sense.
\item Its dominant part is a gamma function, so the first approximation is
      an inverse gamma function rather than merely a logarithm.
\item Inverting Stirling's formula introduces a Lambert $W$ function and a
      polynomial--logarithmic expansion.
\item The Bell correction is absolutely tiny compared with $\Gamma(x+1)$,
      but after inversion it survives as a new hierarchy of sectors of
      order $y^{-n}$, each carrying an oscillatory phase and its own
      inverse-power/logarithmic coefficient expansion.
\end{enumerate}

The linked ProveIt draft collection develops arithmetic, composition, and
reversion in polynomial--logarithmic transseries and analyzes related
near-identity inversions~\cite{ProveItSeries}.  The present problem adds a
different geometry: a
\emph{rapid core} whose logarithmic derivative grows slowly, plus a bounded
or algebraically decaying correction.  The correct small parameter is not
$1/x$ alone.  It is the relative tail $\Res(x)/\Core(x)$, which becomes a
power of $1/y$ after the core has been inverted.

\subsection{The three scales}

Set $Y=\log y$.  The large inverse-gamma variable satisfies roughly
\[
  \InvCore(y)\asymp \frac{Y}{\log Y}.
\]
Its asymptotic description contains the following nested scales.

\begin{description}[leftmargin=2.5em,style=nextline]
\item[Lambert scale.]
The leading object $Y/W(Y/\e)$ resums the repeated logarithms generated by
reversing $x\log x-x$.

\item[Stirling scale.]
Centered Stirling corrections form an odd inverse-power series in
$u_0^{-1}\asymp (\log Y)/Y$, with rational functions of $W(Y/\e)$ as
coefficients.

\item[Bell sectors.]
The subfactorial tail produces sectors proportional to
$y^{-n}=\e^{-nY}$.  Inside the $n$th sector are powers of
$\InvCore(y)^{-1}$, $\log\InvCore(y)^{-1}$, and phases
$\exp(\ii n\pi\InvCore(y))$.
\end{description}

Thus a typical monomial in the final expansion has the shape
\begin{equation}\label{eq:typical-monomial}
 \e^{-nY}\e^{\ii n\sigma\pi\InvCore(\e^Y)}
 \InvCore(\e^Y)^{-p}
 \bigl(\log\InvCore(\e^Y)\bigr)^{-q},
 \qquad n,p,q\ge0.
\end{equation}
The outer index $n$ is the decisive valuation.  At each fixed $n$, the
remaining coefficient admits an ordinary slow asymptotic expansion.

\subsection{Why the core is kept exact for as long as possible}

The inverse gamma function has its own divergent Stirling expansion.  One
can refine it hyperasymptotically through Binet-type remainder formulae and
exponentially improved gamma expansions~\cite{DLMFGamma,Olver1997}.  Those
refinements need not be mixed prematurely with the Bell sectors.  Our master
formula treats $\InvCore(y)$ exactly and expands only the perturbation around
it.  This ``sector-exact normal form'' has three advantages:

\begin{itemize}
\item it is independent of how the inverse gamma core is evaluated;
\item it preserves the rigorous ordering by powers of $y^{-1}$;
\item any preferred algebraic, convergent, numerical, or hyperasymptotic
      representation of $\InvCore$ can be substituted afterward.
\end{itemize}

\subsection{Main contributions}

The main results may be summarized as follows.

\begin{enumerate}
\item We derive an exact decomposition of every lateral incomplete-gamma
      continuation into a gamma core and an endpoint-Mellin tail.
\item We prove the Bell-number expansion with an explicit signed remainder,
      including derivative expansions needed by reversion.
\item We obtain a centered Lambert--Stirling expansion of the large inverse
      gamma branch to arbitrary order, with a triangular coefficient
      recurrence and explicit initial terms.
\item We prove an all-orders Lagrange--B\"urmann formula for inversion around
      a rapid core and rewrite it as a valuation-raising differential
      operator calculus.
\item We specialize that calculus to the principal and conjugate
      subfactorial continuations, obtaining explicit $y^{-1}$ and $y^{-2}$
      sectors and a closed coefficient formula for every sector.
\item We analyze the real median interpolation and the discrete generalized
      inverse separately, including exponentially small threshold shifts.
\item We generalize the Bell mechanism to arbitrary endpoint-Mellin tails
      and explain which parts of the inverse transseries are universal.
\end{enumerate}

\section{Subfactorials beyond the integers}
\label{sec:continuations}

\subsection{Elementary integer formulae}

The inclusion--exclusion formula \eqref{eq:Dn-def} gives the familiar
recurrences
\begin{equation}\label{eq:derangement-recurrences}
 D_n=nD_{n-1}+(-1)^n
     =(n-1)(D_{n-1}+D_{n-2}),
 \qquad D_0=1,\quad D_1=0.
\end{equation}
The first form already hints at a factorial core plus a parity-dependent
small correction.  To expose that correction exactly, write
\begin{equation}\label{eq:tail-of-e}
 \e^{-1}-\sum_{j=0}^{n}\frac{(-1)^j}{j!}
 =\sum_{j=n+1}^{\infty}\frac{(-1)^j}{j!}.
\end{equation}
Repeated integration by parts or Taylor's theorem with integral remainder
will convert this tail into the endpoint integral introduced below.

\subsection{Lateral incomplete-gamma continuations}

The upper incomplete gamma function at a negative argument depends on the
chosen determination of the logarithm in $t^{z}$.  We use the following
notation.

\begin{definition}[Lateral subfactorials]\label{def:lateral-sub}
For $\sigma\in\{+1,-1\}$, let
\begin{equation}\label{eq:lateral-def}
 \Sub_{\sigma}(z)=\frac1\e\Gamma_{\sigma}(z+1,-1),
\end{equation}
where the integration path reaches the negative real axis with
$\arg(-1)=\sigma\pi$.  The choice $\sigma=+1$ is the usual principal
boundary value.  More generally, after winding $k\in\mathbb Z$ times around
the origin one obtains a continuation with phase
$\omega_k=(2k+1)\pi$.
\end{definition}

Define the two building blocks
\begin{equation}\label{eq:A-J-def}
 \Core(z)=\frac{\Gamma(z+1)}\e,
 \qquad
 \Tail(z)=\frac1\e\int_0^1\e^t t^z\dd t,
 \qquad \Re z>-1.
\end{equation}

\begin{theorem}[Exact core--tail decomposition]\label{thm:core-tail}
For $\Re z>-1$,
\begin{equation}\label{eq:core-tail}
 \boxed{\quad
 \Sub_{\sigma}(z)=\Core(z)+\e^{\ii\sigma\pi z}\Tail(z).
 \quad}
\end{equation}
More generally, the branch with $\arg(-1)=(2k+1)\pi$ is
\begin{equation}\label{eq:winding-branch}
 \Sub_k(z)=\Core(z)+\e^{\ii(2k+1)\pi z}\Tail(z).
\end{equation}
Every one of these functions continues to an entire function of $z$, and
all of them agree with $D_n$ at nonnegative integers.
\end{theorem}

\begin{proof}
For $s=z+1$ and the upper lateral path, parameterize the segment from $0$
to $-1$ by $t=-u$ with $\arg t=\pi$.  Then
\[
 \gamma_+(s,-1)=\int_0^{-1}t^{s-1}\e^{-t}\dd t
 =\e^{\ii\pi s}\int_0^1u^{s-1}\e^u\dd u.
\]
The relation $\Gamma_+(s,-1)=\Gamma(s)-\gamma_+(s,-1)$ and
$\e^{\ii\pi(z+1)}=-\e^{\ii\pi z}$ yield
\[
 \Gamma_+(z+1,-1)=\Gamma(z+1)
 +\e^{\ii\pi z}\int_0^1u^z\e^u\dd u.
\]
The lower path is identical with $\pi$ replaced by $-\pi$, and additional
windings replace it by $(2k+1)\pi$.

For fixed nonzero second argument, the upper incomplete gamma function is
entire in its first parameter.  One can also see cancellation directly from
\eqref{eq:J-meromorphic} below: at $z=-m$, $m\ge1$, the residues of
$\Core(z)$ and $\e^{\ii(2k+1)\pi z}\Tail(z)$ are respectively
$(-1)^{m-1}/(\e(m-1)!)$ and $(-1)^m/(\e(m-1)!)$.  They cancel.  Finally,
at $z=n\in\mathbb N_0$, the phase is $(-1)^n$; comparison with the integral
remainder for the exponential series gives $\Sub_k(n)=D_n$.
\end{proof}

\begin{corollary}[Exact integer correction]\label{cor:integer-correction}
For every $n\ge0$,
\begin{equation}\label{eq:integer-correction}
 D_n=\frac{n!}{\e}+(-1)^n\Tail(n).
\end{equation}
Moreover,
\begin{equation}\label{eq:J-simple-bound}
 0<\Tail(x)<\frac1{x+1},\qquad x>-1.
\end{equation}
Consequently, $D_n$ is the nearest integer to $n!/\e$ for every $n\ge1$.
\end{corollary}

\begin{proof}
The first assertion is \eqref{eq:core-tail} at an integer.  Since
$0<\e^{t-1}<1$ for $0\le t<1$,
\[
 0<\Tail(x)=\int_0^1\e^{t-1}t^x\dd t
 <\int_0^1t^x\dd t=\frac1{x+1}.
\]
For $n\ge1$ this is strictly less than $1/2$, proving the nearest-integer
statement.
\end{proof}

\subsection{A canonical real interpolation}

The lateral values are conjugate on the real axis.  Their arithmetic mean
is therefore a natural real entire interpolation.

\begin{definition}[Median subfactorial]\label{def:median}
Define
\begin{equation}\label{eq:median-def}
 \Sub_{\med}(z)=\frac{\Sub_+(z)+\Sub_-(z)}2
 =\Core(z)+\cos(\pi z)\Tail(z).
\end{equation}
\end{definition}

The residue cancellation in the proof of \cref{thm:core-tail} shows that
$\Sub_{\med}$ is entire.  It is real-valued on the real axis and agrees with
$D_n$ at every nonnegative integer.  It is not the only possible real
interpolation; without a normalization condition, infinitely many entire
functions can be added while vanishing on the integers.  It is, however, the
canonical median of the two nearest lateral determinations, and its inverse
will provide a useful real counterpart to the principal complex inverse.

\begin{proposition}[Eventual monotonicity]\label{prop:eventual-mono}
The real function $\Sub_{\med}(x)$ is strictly increasing for all sufficiently
large $x$.  The same is true of the modulus-dominant gamma core
$\Core(x)$, and
\begin{equation}\label{eq:median-derivative}
 \Sub_{\med}'(x)
 =\Core(x)\psi(x+1)
 +\Tail'(x)\cos(\pi x)-\pi\Tail(x)\sin(\pi x).
\end{equation}
\end{proposition}

\begin{proof}
As $x\to\infty$, $\Core(x)\to\infty$ and
$\psi(x+1)\sim\log x$.  On the other hand,
$\Tail(x)=\Ord(x^{-1})$ and $\Tail'(x)=\Ord(x^{-2})$, as follows either
from the integral or from \cref{thm:J-Bell}.  Hence the first term in
\eqref{eq:median-derivative} dominates the other two and is positive.
\end{proof}

\section{The endpoint tail and its Bell-number expansion}
\label{sec:Bell-tail}

The small function $\Tail$ contains the entire derangement correction.  Its
asymptotics are especially transparent because it is simultaneously an
endpoint integral, a meromorphic partial-fraction series, and a Poisson
moment transform.

\subsection{Exact series and meromorphic continuation}

\begin{proposition}[Partial-fraction representation]\label{prop:J-series}
For $\Re z>-1$,
\begin{equation}\label{eq:J-meromorphic}
 \Tail(z)=\frac1\e\sum_{j=0}^{\infty}
 \frac1{j!(z+j+1)}.
\end{equation}
The series converges locally uniformly away from
$z=-1,-2,\ldots$ and gives the meromorphic continuation of $\Tail$ to
$\mathbb C$, with simple residues
\begin{equation}\label{eq:J-residues}
 \Resi_{z=-m}\Tail(z)=\frac1{\e(m-1)!},\qquad m\ge1.
\end{equation}
For every nonnegative integer $r$,
\begin{equation}\label{eq:J-derivative-series}
 \Tail^{(r)}(z)=\frac{(-1)^r r!}{\e}
 \sum_{j=0}^{\infty}\frac1{j!(z+j+1)^{r+1}}.
\end{equation}
\end{proposition}

\begin{proof}
Expand $\e^t=\sum_{j\ge0}t^j/j!$ in \eqref{eq:A-J-def} and integrate
termwise.  Uniform convergence on $[0,1]$ and local uniform convergence in
$z$ on compact subsets of $\Re z>-1$ justify the interchange.  The resulting
series converges normally on compact subsets avoiding its displayed poles,
which proves the continuation and residues.  Termwise differentiation gives
\eqref{eq:J-derivative-series}.
\end{proof}

\subsection{Dobi\'nski moments shifted by one}

Recall Dobi\'nski's formula
\begin{equation}\label{eq:Dobinski}
 \Bell_k=\frac1\e\sum_{j=0}^{\infty}\frac{j^k}{j!},
 \qquad k\ge0.
\end{equation}
A shifted version is the form naturally required by \eqref{eq:J-meromorphic}.

\begin{lemma}[Shifted Poisson moment]\label{lem:shifted-Dobinski}
For every integer $r\ge0$,
\begin{equation}\label{eq:shifted-Dobinski}
 \frac1\e\sum_{j=0}^{\infty}\frac{(j+1)^r}{j!}=\Bell_{r+1}.
\end{equation}
\end{lemma}

\begin{proof}
Let $N$ be Poisson with mean $1$.  Then the left-hand side is
$\mathbb E(N+1)^r$.  Expanding by the binomial theorem and using
$\mathbb E N^q=\Bell_q$ gives
\[
 \mathbb E(N+1)^r=\sum_{q=0}^{r}\binom rq\Bell_q.
\]
The Bell recurrence
$\Bell_{r+1}=\sum_{q=0}^{r}\binom rq\Bell_q$ completes the proof.  Equivalently,
one may compare coefficients in
$\e^{\e^t-1+t}=\frac{\dd}{\dd t}\e^{\e^t-1}$.
\end{proof}

\subsection{The Bell expansion with an exact remainder}

\begin{theorem}[Bell expansion and sharp elementary bound]
\label{thm:J-Bell}
Let $m\ge1$ be fixed and $x>0$.  Then
\begin{equation}\label{eq:J-Bell}
 \Tail(x)=\sum_{k=1}^{m}(-1)^{k-1}\frac{\Bell_k}{x^k}
 +R_m(x),
\end{equation}
where the remainder has the exact representation
\begin{equation}\label{eq:J-Bell-rem}
 R_m(x)=\frac{(-1)^m}{\e x^m}
 \sum_{j=0}^{\infty}
 \frac{(j+1)^m}{j!(x+j+1)}.
\end{equation}
In particular,
\begin{equation}\label{eq:J-Bell-sign-bound}
 (-1)^mR_m(x)>0,
 \qquad
 |R_m(x)|\le \frac{\Bell_{m+1}}{x^{m+1}}.
\end{equation}
Consequently,
\begin{equation}\label{eq:J-Bell-asymptotic}
 \Tail(x)\sim
 \frac1x-\frac2{x^2}+\frac5{x^3}-\frac{15}{x^4}
 +\frac{52}{x^5}-\frac{203}{x^6}
 +\frac{877}{x^7}-\frac{4140}{x^8}+\cdots.
\end{equation}
\end{theorem}

\begin{proof}
For $a>0$ the finite geometric identity gives
\begin{equation}\label{eq:reciprocal-finite}
 \frac1{x+a}=\sum_{k=1}^{m}(-1)^{k-1}\frac{a^{k-1}}{x^k}
 +(-1)^m\frac{a^m}{x^m(x+a)}.
\end{equation}
Insert $a=j+1$ into \eqref{eq:J-meromorphic} and sum over $j$.  By
\cref{lem:shifted-Dobinski}, the coefficient of $x^{-k}$ is
$(-1)^{k-1}\Bell_k$, and the remaining term is
\eqref{eq:J-Bell-rem}.  Every summand in the latter is positive after
multiplication by $(-1)^m$.  Since $(x+j+1)^{-1}\le x^{-1}$,
\[
 |R_m(x)|\le\frac1{\e x^{m+1}}
 \sum_{j=0}^{\infty}\frac{(j+1)^m}{j!}
 =\frac{\Bell_{m+1}}{x^{m+1}}.
\]
The displayed coefficients are the first Bell numbers.
\end{proof}

\begin{corollary}[Hassani's derangement expansion]\label{cor:Hassani}
For every fixed $m\ge1$ and every integer $n\ge1$,
\begin{equation}\label{eq:Hassani-derived}
 D_n=\frac{n!}{\e}
 +\sum_{k=1}^{m}(-1)^{n+k-1}\frac{\Bell_k}{n^k}
 +(-1)^nR_m(n),
\end{equation}
where $|R_m(n)|\le \Bell_{m+1}n^{-m-1}$.
\end{corollary}

\begin{proof}
Substitute \eqref{eq:J-Bell} into \eqref{eq:integer-correction}.
\end{proof}

The proof shows slightly more than a Poincar\'e expansion: successive
truncations alternately bound $\Tail(x)$.  This sign information is useful
for certified approximations.

\subsection{Derivative expansions}

The inverse coefficients involve derivatives of $\Tail$.  They can be
obtained by differentiating the asymptotic expansion, but it is useful to
record the result explicitly.

\begin{proposition}[Differentiated Bell expansion]\label{prop:J-derivatives}
For fixed integers $r\ge0$ and $m\ge1$,
\begin{equation}\label{eq:J-derivative-Bell}
 \Tail^{(r)}(x)
 =\sum_{k=1}^{m}(-1)^{r+k-1}(k)_r\Bell_k\,x^{-k-r}
 +\Ord_{m,r}(x^{-m-r-1}),
\end{equation}
where $(k)_r=k(k+1)\cdots(k+r-1)$ is the rising factorial.  In particular,
\begin{align}
 \Tail'(x)&=-x^{-2}+4x^{-3}-15x^{-4}+60x^{-5}
            -260x^{-6}+\cdots,\label{eq:Jprime-first}\\
 \Tail''(x)&=2x^{-3}-12x^{-4}+60x^{-5}-300x^{-6}
            +1560x^{-7}+\cdots.\label{eq:Jsecond-first}
\end{align}
\end{proposition}

\begin{proof}
Apply \eqref{eq:J-derivative-series} and expand
$(x+j+1)^{-r-1}$ in powers of $x^{-1}$, or differentiate
\eqref{eq:J-Bell} with a remainder obtained from
\eqref{eq:J-Bell-rem}.  For fixed $r,m$, the differentiated remainder is
bounded by a convergent Poisson moment and has the asserted order.
\end{proof}

\subsection{Watson's lemma and the origin of the Bell numbers}

There is a conceptual explanation that generalizes immediately.  With
$t=\e^{-s}$,
\begin{equation}\label{eq:J-Laplace}
 \Tail(x)=\int_0^{\infty}\e^{-(x+1)s}
 \exp(\e^{-s}-1)\dd s.
\end{equation}
The Bell exponential generating function gives
\begin{equation}\label{eq:Bell-egf-minus}
 \exp(\e^{-s}-1)=\sum_{r=0}^{\infty}\Bell_r\frac{(-s)^r}{r!}.
\end{equation}
Watson's lemma therefore yields the equivalent shifted expansion
\begin{equation}\label{eq:J-shifted-Bell}
 \Tail(x)\sim\sum_{r=0}^{\infty}
 \frac{(-1)^r\Bell_r}{(x+1)^{r+1}}.
\end{equation}
Re-expanding $(x+1)^{-r-1}$ in $x^{-1}$ converts
\eqref{eq:J-shifted-Bell} into \eqref{eq:J-Bell-asymptotic}.  The shift by
one in \cref{lem:shifted-Dobinski} is exactly the algebraic shadow of the
Jacobian $\dd t=-\e^{-s}\dd s$.

\begin{summarybox}
The Bell numbers enter before inversion.  They are the endpoint derivatives
of $\exp(\e^{-s}-1)$, equivalently the shifted Poisson moments of the
partial fractions in \eqref{eq:J-meromorphic}.  In the inverse problem they
remain inside the slowly varying coefficient functions of the $y^{-n}$
sectors; inversion does not replace them by a different combinatorial
sequence, but combines them through products, derivatives, and Bell
polynomials of a second kind.
\end{summarybox}

\section{The inverse-gamma core}
\label{sec:inverse-gamma}

Before the Bell perturbation can be inverted, the factorial core must be
understood.  Let $x_\ast>0$ denote the unique positive zero of
$\psi(x+1)$; then $\Core$ is strictly increasing on $(x_\ast,\infty)$.
For every sufficiently large $y$ there is a unique real number
$\InvCore(y)>x_\ast$ satisfying
\begin{equation}\label{eq:a-def}
 \Core(\InvCore(y))=\frac{\Gamma(\InvCore(y)+1)}\e=y.
\end{equation}
In terms of the principal real inverse gamma branch,
\begin{equation}\label{eq:a-invgamma}
 \InvCore(y)=\Gamma_0^{-1}(\e y)-1.
\end{equation}
The theory of inverse gamma branches has been studied analytically and
computationally by Uchiyama, Pedersen, Borwein, Corless, Jeffrey, and their
collaborators~\cite{Uchiyama2012,Pedersen2015,BorweinCorless2018,
CorlessAmenyouJeffrey2017}.

\subsection{Exact differential identities}

Set
\begin{equation}\label{eq:P-def}
 P(x)=\psi(x+1)=\frac{\Core'(x)}{\Core(x)}.
\end{equation}
Implicit differentiation of \eqref{eq:a-def} gives
\begin{equation}\label{eq:a-prime}
 \InvCore'(y)=\frac1{yP(\InvCore(y))}.
\end{equation}
Further derivatives are obtained recursively.  For example, writing
$a=\InvCore(y)$,
\begin{align}
 \InvCore''(y)
 &=-\frac1{y^2P(a)}-\frac{P'(a)}{y^2P(a)^3},
 \label{eq:a-second}\\
 \InvCore'''(y)
 &=\frac{2}{y^3P(a)}
 +\frac{3P'(a)}{y^3P(a)^3}
 -\frac{P''(a)}{y^3P(a)^4}
 +\frac{3P'(a)^2}{y^3P(a)^5}.
 \label{eq:a-third}
\end{align}
These identities are useful both for error propagation and for direct
Lagrange--B\"urmann calculations.

\subsection{Centered Stirling coordinates}

The ordinary Stirling expansion becomes particularly sparse after the
half-unit shift
\begin{equation}\label{eq:u-def}
 u=a+\frac12.
\end{equation}
For $u\to+\infty$,
\begin{equation}\label{eq:centered-stirling}
 \log\Gamma\left(u+\frac12\right)-\frac12\log(2\pi)
 \sim u\log u-u+\sum_{m=1}^{\infty}\frac{c_m}{u^{2m-1}},
\end{equation}
where
\begin{equation}\label{eq:c-m}
 c_m=\frac{B_{2m}(1/2)}{2m(2m-1)}
 =\frac{(2^{1-2m}-1)B_{2m}}{2m(2m-1)}.
\end{equation}
The first coefficients are
\begin{equation}\label{eq:c-first}
 c_1=-\frac1{24},\qquad
 c_2=\frac7{2880},\qquad
 c_3=-\frac{31}{40320},\qquad
 c_4=\frac{127}{215040}.
\end{equation}
Since $\Gamma(a+1)=\e y$, define
\begin{equation}\label{eq:lambda-def}
 \lambda=\log\frac{\e y}{\sqrt{2\pi}}.
\end{equation}
Then the exact inversion problem has the asymptotic normal form
\begin{equation}\label{eq:gamma-normal-form}
 u\log u-u+\sum_{m\ge1}c_m u^{-(2m-1)}\sim\lambda.
\end{equation}

\subsection{The Lambert seed}

Ignore the inverse-power corrections in \eqref{eq:gamma-normal-form}.  The
remaining equation
\begin{equation}\label{eq:seed-equation}
 u_0(\log u_0-1)=\lambda
\end{equation}
can be solved exactly by Lambert $W$.  Put
\begin{equation}\label{eq:w-u0-def}
 w=W_0(\lambda/\e),
 \qquad
 u_0=\frac{\lambda}{w}=\e^{w+1}.
\end{equation}
Then $\log u_0-1=w$, and \eqref{eq:seed-equation} follows.  This is the
natural zeroth approximation to the principal inverse gamma branch; it is
the centered form of the inverse-factorial approximation discussed by
Borwein and Corless~\cite{BorweinCorless2018}.

\subsection{All-orders centered inversion}

\begin{theorem}[Centered inverse-gamma expansion]\label{thm:inverse-gamma-expansion}
As $y\to+\infty$, with $\lambda,w,u_0$ defined by
\eqref{eq:lambda-def}--\eqref{eq:w-u0-def},
\begin{equation}\label{eq:a-pm-expansion}
 \InvCore(y)=-\frac12+u_0+
 \sum_{m=1}^{M}\frac{p_m(w)}{u_0^{2m-1}}
 +\Ord_M(u_0^{-2M-1}),
\end{equation}
where each $p_m(w)$ is a rational function regular for $w\ge0$.  The first
four are
\begin{align}
 p_1(w)&=\frac1{24(1+w)},\label{eq:p1}\\
 p_2(w)&=-\frac{14w^2+38w+29}{5760(1+w)^3},\label{eq:p2}\\
 p_3(w)&=\frac{2232w^4+10104w^3+17634w^2+14234w+4577}
 {2903040(1+w)^5},\label{eq:p3}\\
 p_4(w)&=-\frac{1}{1393459200(1+w)^7}
 \Bigl(822960w^6+5230296w^5+13944036w^4\notag\\
 &\hspace{5.5em}{}+20000424w^3+16328586w^2+7230574w+1368337\Bigr).
 \label{eq:p4}
\end{align}
The coefficients are determined uniquely by the triangular recurrence in
\eqref{eq:pm-recurrence} below.
\end{theorem}

\begin{proof}
Write $r=u_0^{-1}$ and seek
\begin{equation}\label{eq:Delta-series}
 u=u_0+\Delta(r),
 \qquad
 \Delta(r)=\sum_{m\ge1}p_m(w)r^{2m-1}.
\end{equation}
Since $\log u_0-1=w$, subtraction of
$u_0\log u_0-u_0=\lambda$ from \eqref{eq:gamma-normal-form} gives
\begin{align}
0\sim\mathcal E(\Delta;r,w):={}&
 w\Delta+(r^{-1}+\Delta)\log(1+r\Delta)\notag\\
&+\sum_{j\ge1}c_jr^{2j-1}(1+r\Delta)^{-(2j-1)}.
\label{eq:E-Delta}
\end{align}
The coefficient of $r^{2m-1}$ is affine in $p_m$, and its coefficient of
$p_m$ is $1+w$.  Therefore the equations can be solved successively.  The
absence of even powers follows inductively from the odd centered Stirling
series and the fact that $r\Delta$ is even.  Direct coefficient extraction
gives \eqref{eq:p1}--\eqref{eq:p4}.  Standard asymptotic implicit-function
arguments, or substitution followed by one Newton correction, prove the
stated remainder at each fixed $M$.
\end{proof}

For an explicit recurrence, let
\begin{equation}\label{eq:Delta-less-m}
 \Delta_{<m}(r)=\sum_{j=1}^{m-1}p_j(w)r^{2j-1}.
\end{equation}
Then
\begin{align}
 p_m(w)=-\frac1{1+w}\coeff{r}{2m-1}
 \Bigg\{&w\Delta_{<m}
 +(r^{-1}+\Delta_{<m})\log(1+r\Delta_{<m})\notag\\
 &+\sum_{j=1}^{m}c_jr^{2j-1}
  (1+r\Delta_{<m})^{-(2j-1)}\Bigg\}.
 \label{eq:pm-recurrence}
\end{align}
This recurrence uses only addition, multiplication, rational powers of a
unit series, and coefficient extraction.  It is therefore directly
implementable in any formal power-series system.

\subsection{Expansion into iterated logarithms}

The Lambert seed can itself be expanded.  Let
\begin{equation}\label{eq:L1L2-def}
 L_1=\log(\lambda/\e),
 \qquad L_2=\log L_1.
\end{equation}
The standard large-$z$ expansion of $W_0(z)$~\cite{CorlessEtAl1996} implies
\begin{align}
 \frac1{W_0(\lambda/\e)}={}&
 \frac1{L_1}+\frac{L_2}{L_1^2}
 +\frac{L_2^2-L_2}{L_1^3}
 +\frac{L_2^3-\frac52L_2^2+L_2}{L_1^4}\notag\\
 &+\frac{L_2^4-\frac{13}{3}L_2^3+\frac92L_2^2-L_2}{L_1^5}
 +\Ord\!\left(\frac{L_2^5}{L_1^6}\right).
 \label{eq:inverse-W-expansion}
\end{align}
Consequently,
\begin{equation}\label{eq:u0-log-expansion}
 u_0=\lambda\left[
 \frac1{L_1}+\frac{L_2}{L_1^2}
 +\frac{L_2^2-L_2}{L_1^3}+\cdots\right].
\end{equation}
Substituting \eqref{eq:u0-log-expansion} into
\eqref{eq:a-pm-expansion} produces a purely polynomial--logarithmic
expansion in $\lambda^{-1}$, $L_1$, and $L_2$.  Keeping $W_0$ unexpanded is
usually more compact and numerically superior.

\subsection{Residual certification and Newton improvement}

Let $a_M(y)$ denote the truncation in \eqref{eq:a-pm-expansion}, and define
the logarithmic residual
\begin{equation}\label{eq:gamma-residual}
 \rho_M(y)=\log\Gamma(a_M(y)+1)-1-\log y.
\end{equation}
Since the derivative with respect to $a$ is $P(a)$, the first omitted error
is certified by
\begin{equation}\label{eq:gamma-error-residual}
 \InvCore(y)-a_M(y)
 =-\frac{\rho_M(y)}{P(a_M(y))}
 +\Ord\!\left(\frac{\rho_M(y)^2P'(a_M(y))}{P(a_M(y))^3}\right).
\end{equation}
One Newton step,
\begin{equation}\label{eq:gamma-Newton}
 a_M^{\mathrm N}=a_M-\frac{\rho_M}{P(a_M)},
\end{equation}
therefore squares the local residual.  This observation will recur for the
full subfactorial inverse.

\begin{checkpointbox}
The centered expansion is already effective at modest arguments.  For
$y=10^6$, the Lambert seed errs by about $1.72\times10^{-3}$; successively
including $p_1,p_2,p_3,p_4$ reduces the absolute error to approximately
$1.27\times10^{-6}$, $3.52\times10^{-9}$,
$2.29\times10^{-11}$, and $2.85\times10^{-13}$, respectively.  A detailed
table appears in \cref{sec:numerics}.
\end{checkpointbox}

\section{A general reversion calculus for a rapid core plus a small tail}
\label{sec:general-calculus}

We now isolate the mechanism that will be applied to the subfactorial.  The
basic problem is
\begin{equation}\label{eq:general-F}
 F(x)=\Core(x)+\Res(x),
 \qquad y=F(x),
\end{equation}
where $\Core$ is eventually one-to-one and grows so rapidly that
$\Res(x)/\Core(x)\to0$.  The inverse of $\Core$ will be denoted by
\begin{equation}\label{eq:general-a}
 a(y)=\Core^{-1}(y).
\end{equation}

\subsection{Lagrange--B\"urmann after straightening the core}

The most efficient derivation begins by using the core itself as a
coordinate.  Put
\begin{equation}\label{eq:rho-def}
 u=\Core(x),
 \qquad x=a(u),
 \qquad \rho(u)=\Res(a(u)).
\end{equation}
Then \eqref{eq:general-F} becomes the near-identity equation
\begin{equation}\label{eq:near-identity-u}
 y=u+\rho(u).
\end{equation}
Ordinary Lagrange--B\"urmann inversion can now be used without attempting to
reverse the rapid function directly.

\begin{theorem}[Master reversion identity]\label{thm:master-reversion}
Introduce a bookkeeping parameter $\varepsilon$ and let $x_\varepsilon(y)$
be the local inverse of
\begin{equation}\label{eq:F-epsilon}
 F_\varepsilon(x)=\Core(x)+\varepsilon\Res(x)
\end{equation}
that satisfies $x_0(y)=a(y)$.  Whenever the local analytic inverse exists,
\begin{equation}\label{eq:master-LB}
 \boxed{\quad
 x_\varepsilon(y)=a(y)+
 \sum_{n=1}^{\infty}\frac{(-\varepsilon)^n}{n!}
 \frac{\dd^{n-1}}{\dd y^{n-1}}
 \left[a'(y)\Res(a(y))^n\right].
 \quad}
\end{equation}
The identity is valid as a convergent Taylor series in $\varepsilon$ near
zero and, independently, as a formal identity.  If the coefficients are
asymptotically ordered when $y\to\infty$, setting $\varepsilon=1$ gives the
inverse transseries of $F$.
\end{theorem}

\begin{proof}
In the straightened coordinate $u$, the equation is
$y=u+\varepsilon\rho(u)$.  The Lagrange--B\"urmann formula for an analytic
test function $H$ states
\begin{equation}\label{eq:LB-H}
 H(u(y))=H(y)+\sum_{n=1}^{\infty}
 \frac{(-\varepsilon)^n}{n!}
 \frac{\dd^{n-1}}{\dd y^{n-1}}
 \left[H'(y)\rho(y)^n\right].
\end{equation}
Choose $H=a$.  Since $x=a(u)$ and $\rho(y)=\Res(a(y))$, this is exactly
\eqref{eq:master-LB}.
\end{proof}

The first three terms, written without any structural assumptions, are
\begin{align}
 x(y)={}&a-a'R
 +\frac12\frac{\dd}{\dd y}(a'R^2)
 -\frac16\frac{\dd^2}{\dd y^2}(a'R^3)+\cdots,
 \label{eq:master-first-three}
\end{align}
where $R=\Res(a(y))$.  The universal leading displacement is therefore
\begin{equation}\label{eq:universal-leading-shift}
 x-a\sim-\frac{\Res(a)}{\Core'(a)}.
\end{equation}
This is the familiar ``residual divided by slope'' rule, now embedded in an
all-orders formula.

\subsection{The logarithmic derivative and outer valuation}

Define the logarithmic derivative of the core by
\begin{equation}\label{eq:Lambda-general}
 \Lam(x)=\frac{\Core'(x)}{\Core(x)}.
\end{equation}
At $x=a(y)$,
\begin{equation}\label{eq:a-prime-general}
 a'(y)=\frac1{y\Lam(a(y))}.
\end{equation}
For any differentiable coefficient $f$ and integer $r\ge1$,
\begin{equation}\label{eq:valuation-derivative}
 \frac{\dd}{\dd y}\left[y^{-r}f(a(y))\right]
 =y^{-r-1}\mathscr D_r f(a(y)),
\end{equation}
where
\begin{equation}\label{eq:D-r-general}
 \mathscr D_r f=\frac{f'}{\Lam}-rf.
\end{equation}
Every $y$-derivative therefore raises the outer $y^{-1}$ valuation by one.
This simple fact is the source of local finiteness in the resulting
transseries.

\begin{theorem}[Operator form of the inverse sectors]
\label{thm:operator-sectors}
Assume that $\Core$ is eventually positive and increasing, and let
$\Lam=\Core'/\Core$.  Then the formal inverse of
$F=\Core+\Res$ around $a=\Core^{-1}$ is
\begin{equation}\label{eq:general-sector-expansion}
 F^{-1}(y)=a(y)+\sum_{n=1}^{\infty}\frac{X_n(a(y))}{y^n},
\end{equation}
with
\begin{equation}\label{eq:Xn-general}
 \boxed{\quad
 X_n(x)=\frac{(-1)^n}{n!}
 \bigl(\mathscr D_{n-1}\circ\cdots\circ\mathscr D_1\bigr)
 \left(\frac{\Res(x)^n}{\Lam(x)}\right),
 \quad}
\end{equation}
where an empty composition is the identity.  In particular,
\begin{align}
 X_1&=-\frac{\Res}{\Lam},\label{eq:X1-general}\\
 X_2&=\frac12\left[
 \frac1\Lam\left(\frac{\Res^2}{\Lam}\right)'
 -\frac{\Res^2}{\Lam}\right]\label{eq:X2-general-a}\\
 &=\frac{\Res\Res'}{\Lam^2}
 -\frac{\Core''}{2\Core}\frac{\Res^2}{\Lam^3}.
 \label{eq:X2-general-b}
\end{align}
\end{theorem}

\begin{proof}
In the $n$th term of \eqref{eq:master-LB}, use
$a'=1/(y\Lam(a))$ to obtain
\[
 a'(y)\Res(a(y))^n
 =y^{-1}\left(\frac{\Res^n}{\Lam}\right)_{x=a(y)}.
\]
Applying \eqref{eq:valuation-derivative} successively for
$r=1,2,\ldots,n-1$ proves \eqref{eq:Xn-general}.  The first two
coefficients follow directly.  To pass from \eqref{eq:X2-general-a} to
\eqref{eq:X2-general-b}, note that
$\Lam'=\Core''/\Core-\Lam^2$; the $\Res^2/(2\Lam)$ terms cancel.
\end{proof}

\subsection{An admissible asymptotic class and remainder theorem}

The operator identity is formal without any estimates.  The following
conditions are convenient for a Poincar\'e interpretation.  They are not
minimal, but they cover the subfactorial and many gamma-type examples.

\begin{definition}[Admissible rapid-core pair]\label{def:admissible-pair}
A pair $(\Core,\Res)$ on $[x_0,\infty)$ is called admissible if:
\begin{enumerate}
\item $\Core>0$ is analytic in a fixed complex neighborhood of the ray,
      is strictly increasing there, and $\Core(x)\to\infty$;
\item $\Lam=\Core'/\Core>0$ and there is a positive scale $s(x)$ with
      $s(x)=\littleo(\Lam(x))$;
\item for every fixed $k\ge0$,
\[
 \Res^{(k)}(x)=\Ord(M(x)s(x)^k),
 \qquad
 (\Lam^{\pm1})^{(k)}(x)
 =\Ord(\Lam(x)^{\pm1}s(x)^k),
\]
for some positive $M(x)$;
\item $M(x)/\Core(x)\to0$.
\end{enumerate}
Oscillatory factors $\e^{\ii\omega x}$ with fixed $\omega$ are permitted by
replacing $s$ with $s+|\omega|$.
\end{definition}

\begin{theorem}[Fixed-sector asymptotic remainder]
\label{thm:general-remainder}
For an admissible pair, suppose in addition that the local inverse branch of
$\Core+\varepsilon\Res$ remains analytic for $|\varepsilon|\le1+\eta$ for
all sufficiently large $y$, with some fixed $\eta>0$.  Then, for every fixed
$N$,
\begin{equation}\label{eq:general-remainder}
 F^{-1}(y)=a(y)+\sum_{n=1}^{N}\frac{X_n(a(y))}{y^n}
 +\Ord_N\!\left(
 \frac{M(a(y))^{N+1}}{y^{N+1}\Lam(a(y))}
 \right).
\end{equation}
Moreover,
\begin{equation}\label{eq:Xn-size}
 X_n(x)=\Ord_n\!\left(\frac{M(x)^n}{\Lam(x)}\right).
\end{equation}
\end{theorem}

\begin{proof}
Because $s/\Lam\to0$, the operator
$\mathscr D_r=\Lam^{-1}D_x-r$ maps a symbol of size $M^n/\Lam$ to one of the
same size at fixed $r,n$.  This proves \eqref{eq:Xn-size} inductively.
Analyticity in $\varepsilon$ and Cauchy's estimate bound the Taylor remainder
of \eqref{eq:master-LB} by the first neglected asymptotic scale, uniformly
for sufficiently large $y$.  The factor $y^{-N-1}$ is explicit from the
valuation-raising identity \eqref{eq:valuation-derivative}; the remaining
coefficient has the size in \eqref{eq:Xn-size}.
\end{proof}

\begin{remark}
The uniform analyticity hypothesis can be replaced by a direct
Newton--Kantorovich or Rouch\'e argument.  In the subfactorial case the slope
of the core is of order $y\log a$, whereas the tail and all of its fixed
derivatives are only algebraic in $a$.  A disk of radius
$C/(ya\log a)$ around $a$ is therefore more than sufficient.
\end{remark}

\subsection{Fourier closure and phase-sensitive operators}

Suppose the tail has a finite or asymptotic Fourier decomposition
\begin{equation}\label{eq:R-Fourier}
 \Res(x)=\sum_{\omega\in\Omega}\e^{\ii\omega x}r_\omega(x),
\end{equation}
where $r_\omega$ are slow symbols.  Products in the $n$th sector create only
frequencies in the $n$-fold additive sumset
\begin{equation}\label{eq:sumset}
 n\Omega=\{\omega_1+\cdots+\omega_n:\omega_j\in\Omega\}.
\end{equation}
The coefficient field is closed because
\begin{equation}\label{eq:phase-operator-general}
 \mathscr D_r\bigl(\e^{\ii\omega x}f(x)\bigr)
 =\e^{\ii\omega x}
 \left[\frac{f'+\ii\omega f}{\Lam}-rf\right].
\end{equation}
Thus inversion never generates arbitrary new frequencies: it generates
additive combinations of frequencies already present in the perturbation.
For a single phase $\omega$, sector $n$ carries exactly the phase
$\e^{\ii n\omega a(y)}$.

\subsection{The frozen-coefficient logarithm}

When $D_x/\Lam$ is asymptotically small on the coefficient field, each
operator $\mathscr D_r$ is dominated by multiplication by $-r$.  Therefore
\begin{equation}\label{eq:Xn-leading-general}
 X_n(x)\sim-\frac1n\frac{\Res(x)^n}{\Lam(x)}.
\end{equation}
Summing these leading pieces gives the useful resummation
\begin{equation}\label{eq:frozen-log}
 F^{-1}(y)\approx a(y)+\frac1{\Lam(a(y))}
 \log\left(1-\frac{\Res(a(y))}{y}\right).
\end{equation}
It should not be mistaken for an exact identity: derivatives of $\Res$ and
$\Lam$ supply the systematic corrections.  Nevertheless,
\eqref{eq:frozen-log} captures the leading term in every outer sector at
once.

\subsection{Residual certification and Newton doubling}

Let
\begin{equation}\label{eq:GN-general}
 G_N(y)=a(y)+\sum_{n=1}^{N}\frac{X_n(a(y))}{y^n}
\end{equation}
be a truncated inverse.  Its residual is
\begin{equation}\label{eq:inverse-residual}
 \mathcal E_N(y)=F(G_N(y))-y.
\end{equation}
Taylor's theorem gives
\begin{equation}\label{eq:residual-to-error}
 F^{-1}(y)-G_N(y)
 =-\frac{\mathcal E_N(y)}{F'(G_N(y))}
 +\Ord\!\left(
 \frac{F''(G_N(y))\mathcal E_N(y)^2}{F'(G_N(y))^3}
 \right).
\end{equation}
Because $F'(a)\sim y\Lam(a)$, a residual of outer order $y^{-N}$ produces
an inverse error of outer order $y^{-N-1}$.  The Newton update
\begin{equation}\label{eq:Newton-general}
 G_N^{\mathrm N}=G_N-\frac{F(G_N)-y}{F'(G_N)}
\end{equation}
then approximately doubles the resolved outer valuation.  This is useful
both numerically and as a symbolic consistency check.

\subsection{A coefficient algorithm}

\begin{algorithm}[Rapid-core inverse sectors]\label{alg:rapid-core}
Given symbolic descriptions of $\Lam(x)$ and $\Res(x)$ and a desired sector
order $N$:
\begin{enumerate}
\item Set $X_0(x)=x$ only as a bookkeeping label.
\item For $n=1,\ldots,N$, initialize
\[
 H\leftarrow\frac{\Res(x)^n}{\Lam(x)}.
\]
\item For $r=1,\ldots,n-1$, replace
\[
 H\leftarrow\frac{H'(x)}{\Lam(x)}-rH.
\]
\item Set $X_n(x)\leftarrow(-1)^nH/n!$.
\item Return
\[
 a(y)+\sum_{n=1}^{N}y^{-n}X_n(a(y)).
\]
\end{enumerate}
If $\Res$ is itself represented by an asymptotic series, truncate its slow
coefficient expansion only after the outer sector $n$ has been fixed.
\end{algorithm}

\begin{warningbox}[title={A scale-ordering trap}]
Expanding $a(y)$ into logarithms, expanding $\Res(a)$ into inverse powers,
and expanding in $y^{-1}$ are three different truncations.  They do not have
the same valuation.  The safe order is: first choose the outer sector order
in $y^{-1}$; next compute its coefficient exactly in $a$; then expand that
coefficient in the slow scale; finally replace $a$ by a Lambert--Stirling
approximation of sufficient accuracy.  Reversing this order can discard
terms that are larger than terms retained elsewhere.
\end{warningbox}

\section{Lateral inverse subfactorials: the complete Bell-sector expansion}
\label{sec:lateral-inverse}

We apply the general calculus to
\begin{equation}\label{eq:sub-as-core-tail-again}
 \Sub_\sigma(x)=\Core(x)+\Res_\sigma(x),
 \qquad
 \Res_\sigma(x)=\e^{\ii\sigma\pi x}\Tail(x),
 \qquad
 \Lam(x)=P(x)=\psi(x+1).
\end{equation}
The core inverse is $a=\InvCore$.

\subsection{Existence and localization of the inverse branches}

\begin{theorem}[Distinguished large inverse branches]
\label{thm:inverse-branch-existence}
There exists $y_0>0$ such that for every real $y>y_0$ and each
$\sigma\in\{+1,-1\}$, the equation
\begin{equation}\label{eq:sub-inverse-equation}
 \Sub_\sigma(z)=y
\end{equation}
has exactly one solution $X_\sigma(y)$ in a disk
\begin{equation}\label{eq:inverse-disk}
 |z-\InvCore(y)|\le
 \frac{C}{y\InvCore(y)\log\InvCore(y)},
\end{equation}
where $C$ is independent of $y$.  The solution is simple, depends
analytically on $y$ in a neighborhood of the positive ray, and satisfies
\begin{equation}\label{eq:branch-leading-localization}
 X_\sigma(y)-\InvCore(y)
 =-\frac{\e^{\ii\sigma\pi\InvCore(y)}\Tail(\InvCore(y))}
 {yP(\InvCore(y))}
 +\Ord\!\left(\frac1{y^2a^2\log a}\right),
\end{equation}
where $a=\InvCore(y)$ in the remainder.
\end{theorem}

\begin{proof}
Let $a=\InvCore(y)$ and write $z=a+\delta$.  On a disk of radius
$r=C/(ya\log a)$, Taylor's theorem and the standard asymptotics of the
polygamma functions give
\[
 \Core(a+\delta)-\Core(a)
 =yP(a)\delta\,[1+\littleo(1)].
\]
On the boundary its modulus is $C/a\,[1+\littleo(1)]$.  Meanwhile
$\e^{\ii\sigma\pi z}\Tail(z)=\Ord(a^{-1})$ uniformly in the disk.  Choosing
$C$ sufficiently large and then $y$ sufficiently large makes the core
difference dominate the tail on the boundary.  Rouch\'e's theorem shows
that the perturbed equation has the same number of zeros as the core
difference, namely one.  Simplicity follows because the derivative is
$yP(a)[1+\littleo(1)]$.  One Taylor step gives
\eqref{eq:branch-leading-localization}; analytic dependence follows from the
implicit-function theorem.
\end{proof}

For real $y$, conjugation of \eqref{eq:core-tail} gives
\begin{equation}\label{eq:conjugate-branches}
 X_-(y)=\overline{X_+(y)}.
\end{equation}
Thus it is enough to compute the principal branch $X_+$, but retaining
$\sigma$ keeps the branch symmetry visible.

\subsection{The phase-adapted differential operators}

For integers $r\ge1$ and $n\ge1$, define
\begin{equation}\label{eq:L-rn-sigma}
 \mathcal L^{(\sigma)}_{r,n}f
 =\frac{f'}{P}
 +\left(\frac{\ii\sigma n\pi}{P}-r\right)f.
\end{equation}
Then
\begin{equation}\label{eq:sub-derivative-rule}
 \frac{\dd}{\dd y}
 \left[y^{-r}\e^{\ii\sigma n\pi a(y)}f(a(y))\right]
 =y^{-r-1}\e^{\ii\sigma n\pi a(y)}
 \mathcal L^{(\sigma)}_{r,n}f(a(y)).
\end{equation}
This is \eqref{eq:phase-operator-general} with $\omega=\sigma n\pi$ and
$\Lam=P$.

\begin{theorem}[All-orders Bell-sector transseries]
\label{thm:sub-all-orders}
For every fixed $N$ and $y\to+\infty$,
\begin{equation}\label{eq:sub-all-orders}
 \boxed{\quad
 X_\sigma(y)=a(y)+\sum_{n=1}^{N}
 \frac{\e^{\ii\sigma n\pi a(y)}}{y^n}C_{n,\sigma}(a(y))
 +\Ord_N\!\left(\frac1{y^{N+1}a(y)^{N+1}\log a(y)}\right),
 \quad}
\end{equation}
where $a=\InvCore$ and
\begin{equation}\label{eq:Cn-sigma}
 \boxed{\quad
 C_{n,\sigma}(x)=\frac{(-1)^n}{n!}
 \bigl(\mathcal L^{(\sigma)}_{n-1,n}\circ\cdots\circ
       \mathcal L^{(\sigma)}_{1,n}\bigr)
 \left(\frac{\Tail(x)^n}{P(x)}\right).
 \quad}
\end{equation}
At fixed $n$,
\begin{equation}\label{eq:Cn-leading}
 C_{n,\sigma}(x)
 \sim-\frac1n\frac{\Tail(x)^n}{P(x)}
 \sim-\frac1{n x^n\log x}.
\end{equation}
\end{theorem}

\begin{proof}
Apply \cref{thm:operator-sectors} to
$\Res_\sigma=\e^{\ii\sigma\pi x}\Tail(x)$ and move the phase
$\e^{\ii\sigma n\pi x}$ outside each phase-adapted operator using
\eqref{eq:sub-derivative-rule}.  Since
$\Tail^{(k)}(x)=\Ord_k(x^{-k-1})$ and
$P(x)\sim\log x$ while $P^{(k)}(x)=\Ord_k(x^{-k})$, the pair is admissible
with $M(x)=x^{-1}$ and a fixed derivative scale.  The remainder follows from
\cref{thm:general-remainder} and the branch existence theorem.

For the leading coefficient, the derivative and phase terms divided by
$P$ are smaller than multiplication by $-r$.  The composition contributes
$(-1)^{n-1}(n-1)!$, which multiplied by $(-1)^n/n!$ gives $-1/n$.
Finally use $\Tail(x)\sim x^{-1}$ and $P(x)\sim\log x$.
\end{proof}

\subsection{The first two sectors}

The first coefficient is particularly simple:
\begin{equation}\label{eq:C1-exact}
 C_{1,\sigma}(x)=-\frac{\Tail(x)}{P(x)}.
\end{equation}
The second is
\begin{equation}\label{eq:C2-exact}
 C_{2,\sigma}(x)=
 \frac{\Tail\Tail'}{P^2}
 +\frac{\ii\sigma\pi\Tail^2}{P^2}
 -\frac{\Tail^2P'}{2P^3}
 -\frac{\Tail^2}{2P}.
\end{equation}
Therefore
\begin{align}
X_\sigma(y)={}&a
 -\frac{\e^{\ii\sigma\pi a}}{y}\frac{\Tail(a)}{P(a)}\notag\\
&+\frac{\e^{2\ii\sigma\pi a}}{y^2}
 \left[
 \frac{\Tail(a)\Tail'(a)}{P(a)^2}
 +\frac{\ii\sigma\pi\Tail(a)^2}{P(a)^2}
 -\frac{\Tail(a)^2P'(a)}{2P(a)^3}
 -\frac{\Tail(a)^2}{2P(a)}
 \right]\notag\\
&+\Ord\!\left(\frac1{y^3a^3\log a}\right).
\label{eq:two-sector-exact}
\end{align}
This form is often more accurate than a much longer fully expanded formula,
because $\Tail$, $P$, and $a$ can all be evaluated directly.

\subsection{Bell--logarithmic expansion inside the first sector}

Let $\ell=\log x$.  The digamma expansion
\begin{equation}\label{eq:P-expansion}
 P(x)=\psi(x+1)
 \sim \ell+\frac1{2x}-\frac1{12x^2}
 +\frac1{120x^4}-\frac1{252x^6}+\cdots
\end{equation}
combined with \eqref{eq:J-Bell-asymptotic} gives
\begin{align}
 \frac{\Tail(x)}{P(x)}\sim{}&
 \frac1{x\ell}
 -\frac{4\ell+1}{2x^2\ell^2}
 +\frac{60\ell^2+13\ell+3}{12x^3\ell^3}\notag\\
&-\frac{360\ell^3+64\ell^2+14\ell+3}{24x^4\ell^4}\notag\\
&+\frac{37440\ell^4+5694\ell^3+1025\ell^2+225\ell+45}
 {720x^5\ell^5}
 +\Ord\!\left(\frac1{x^6\ell}\right).
\label{eq:J-over-P}
\end{align}
The Bell numbers $1,2,5,15,52,\ldots$ are encoded in the numerators after
division by the digamma series.  No new transcendental constants occur in
this first sector.

Substitution into \eqref{eq:C1-exact} yields
\begin{align}
-\frac{\e^{\ii\sigma\pi a}}{y}\frac{\Tail(a)}{P(a)}
={}&-\frac{\e^{\ii\sigma\pi a}}{y}
\Bigg[
 \frac1{a\log a}
 -\frac{4\log a+1}{2a^2(\log a)^2}\notag\\
&+\frac{60(\log a)^2+13\log a+3}
 {12a^3(\log a)^3}
 +\Ord\!\left(\frac1{a^4\log a}\right)
\Bigg].
\label{eq:first-sector-expanded}
\end{align}

\subsection{Expansion inside the second sector}

Expanding \eqref{eq:C2-exact} gives
\begin{align}
 C_{2,\sigma}(x)={}&
 -\frac{\ell-2\ii\sigma\pi}{2x^2\ell^2}\notag\\
&+\frac{8\ell^2-3\ell-2-16\ii\sigma\pi\ell-4\ii\sigma\pi}
 {4x^3\ell^3}\notag\\
&+\frac{-168\ell^3+119\ell^2+75\ell+18
 +336\ii\sigma\pi\ell^2+100\ii\sigma\pi\ell+18\ii\sigma\pi}
 {24x^4\ell^4}\notag\\
&+\Ord\!\left(\frac1{x^5\ell}\right).
\label{eq:C2-expanded}
\end{align}
The appearance of $\pi$ is entirely due to differentiating the oscillatory
phase.  The leading term $-1/(2x^2\ell)$ agrees with the universal law
\eqref{eq:Cn-leading}.

\subsection{A compact nested transseries}

Combining \cref{thm:inverse-gamma-expansion,thm:sub-all-orders} gives a
single all-orders description.  Define
\begin{equation}\label{eq:a-tilde-formal}
 \widehat a(y)=-\frac12+u_0+
 \sum_{m\ge1}\frac{p_m(w)}{u_0^{2m-1}},
 \qquad
 u_0=\frac{\lambda}{W_0(\lambda/\e)},
 \quad \lambda=\log\frac{\e y}{\sqrt{2\pi}}.
\end{equation}
Then, with the understanding that each outer coefficient is expanded only
to the slow order required,
\begin{equation}\label{eq:full-nested-transseries}
 \boxed{\quad
 X_\sigma(y)\sim\widehat a(y)
 +\sum_{n\ge1}\frac{\e^{\ii\sigma n\pi\widehat a(y)}}{y^n}
 \widehat C_{n,\sigma}(\widehat a(y)),
 \quad}
\end{equation}
where $\widehat C_{n,\sigma}$ is obtained from \eqref{eq:Cn-sigma} by
substituting the Bell expansion of $\Tail$ and the Bernoulli expansion of
$P$.  If $W_0$ is further expanded by \eqref{eq:inverse-W-expansion}, every
coefficient becomes a polynomial--logarithmic expression in
$Y=\log y$, $\log Y$, and $\log\log Y$, multiplied by the outer monomial
$\e^{-nY}$ and the chirped phase
$\exp(\ii\sigma n\pi Y/\log Y+\cdots)$.

\subsection{Leading real and imaginary displacements}

For the principal branch $\sigma=+1$, write $a=\InvCore(y)$ and
$Q(a)=\Tail(a)/P(a)>0$.  The first sector gives
\begin{align}
 \Re X_+(y)&=a-\frac{Q(a)}{y}\cos(\pi a)
 +\Ord\!\left(\frac1{y^2a^2\log a}\right),
 \label{eq:principal-real-leading}\\
 \Im X_+(y)&=-\frac{Q(a)}{y}\sin(\pi a)
 +\Ord\!\left(\frac1{y^2a^2\log a}\right).
 \label{eq:principal-imag-leading}
\end{align}
The inverse branch therefore winds across the real axis whenever the core
inverse passes an integer.  At values $y=\Core(n)=n!/\e$, the leading shift
is real because the phase is $(-1)^n$; at generic $y$ it is complex.

\subsection{Leading all-sector resummation}

Applying \eqref{eq:frozen-log} gives
\begin{equation}\label{eq:sub-frozen-log}
 X_\sigma(y)\approx a(y)+\frac1{P(a(y))}
 \log\left(1-\frac{\e^{\ii\sigma\pi a(y)}\Tail(a(y))}{y}\right).
\end{equation}
Expanding the logarithm reproduces
$-\e^{\ii\sigma n\pi a}\Tail(a)^n/[n y^nP(a)]$, the leading term of every
sector.  Formula \eqref{eq:sub-frozen-log} is useful as a compact
multi-sector initial guess for Newton iteration.

\section{The real median inverse}
\label{sec:median-inverse}

The median interpolation
\[
 \Sub_{\med}(x)=\Core(x)+\Tail(x)\cos(\pi x)
\]
is eventually strictly increasing by \cref{prop:eventual-mono}; denote its
large real inverse by $X_{\med}(y)$.  It can be treated either directly with
\cref{thm:operator-sectors} or as the average-frequency perturbation
\[
 \Tail(x)\cos(\pi x)=\frac{\Tail(x)}2
 \left(\e^{\ii\pi x}+\e^{-\ii\pi x}\right).
\]
Unlike the average of the inverse branches, the inverse of the average is a
nonlinear object; higher sectors contain mixed frequencies.

\subsection{All-orders operator formula}

Define
\begin{equation}\label{eq:median-R}
 R_{\med}(x)=\Tail(x)\cos(\pi x),
 \qquad
 \mathscr L_r f=\frac{f'}{P}-rf.
\end{equation}
Then
\begin{theorem}[Median inverse transseries]\label{thm:median-inverse}
For every fixed $N$,
\begin{equation}\label{eq:median-all-orders}
 X_{\med}(y)=a(y)+\sum_{n=1}^{N}\frac{C_n^{\med}(a(y))}{y^n}
 +\Ord_N\!\left(\frac1{y^{N+1}a(y)^{N+1}\log a(y)}\right),
\end{equation}
where
\begin{equation}\label{eq:Cn-median}
 C_n^{\med}(x)=\frac{(-1)^n}{n!}
 \bigl(\mathscr L_{n-1}\circ\cdots\circ\mathscr L_1\bigr)
 \left(\frac{R_{\med}(x)^n}{P(x)}\right).
\end{equation}
The first two sectors are
\begin{align}
 C_1^{\med}(x)&=-\frac{\Tail(x)\cos(\pi x)}{P(x)},
 \label{eq:C1-median}\\
 C_2^{\med}(x)&=\frac12\left[
 \frac1{P(x)}\left(\frac{\Tail(x)^2\cos^2(\pi x)}{P(x)}\right)'
 -\frac{\Tail(x)^2\cos^2(\pi x)}{P(x)}
 \right].
 \label{eq:C2-median}
\end{align}
\end{theorem}

\begin{proof}
This is \cref{thm:operator-sectors,thm:general-remainder} with
$\Res=R_{\med}$ and $\Lam=P$.  The derivative of the cosine is bounded, so
$R_{\med}$ belongs to the same admissible symbol class as the lateral tail.
\end{proof}

The leading approximation is consequently
\begin{equation}\label{eq:median-leading}
 X_{\med}(y)=a(y)
 -\frac{\Tail(a(y))\cos(\pi a(y))}{yP(a(y))}
 +\Ord\!\left(\frac1{y^2a(y)^2\log a(y)}\right).
\end{equation}
Its all-sector frozen-coefficient resummation is
\begin{equation}\label{eq:median-log-resum}
 X_{\med}(y)\approx a(y)+\frac1{P(a(y))}
 \log\left(1-\frac{\Tail(a(y))\cos(\pi a(y))}{y}\right).
\end{equation}

\subsection{Frequency mixing}

At sector $n$, the factor $\cos^n(\pi x)$ contains frequencies
$(n-2j)\pi$, $0\le j\le n$.  Thus
\begin{equation}\label{eq:cos-power}
 \cos^n\theta=2^{-n}\sum_{j=0}^{n}\binom nj
 \e^{\ii(n-2j)\theta}.
\end{equation}
The phase-adapted operators preserve this finite set.  In particular, even
sectors contain a nonoscillatory component, while odd sectors do not.  This
is a useful illustration of the additive-frequency principle in
\cref{eq:sumset}: taking a real median does not remove the oscillatory
algebra; it symmetrizes it.

\subsection{Comparison with the lateral branches}

At first order,
\begin{equation}\label{eq:median-average-first}
 X_{\med}(y)-a(y)
 =\frac{(X_+(y)-a(y))+(X_-(y)-a(y))}{2}
 +\Ord\!\left(\frac1{y^2a^2\log a}\right).
\end{equation}
At second and higher orders this equality fails because inversion is
nonlinear.  The discrepancy can be computed directly from
\eqref{eq:C2-exact} and \eqref{eq:C2-median}; it is generated by mixed
products of the $+\pi$ and $-\pi$ frequencies.

\section{The discrete generalized inverse and its threshold transseries}
\label{sec:discrete-inverse}

The integer sequence itself has the monotone tail
$D_2=1,D_3=2,D_4=9,\ldots$.  Define its right-continuous generalized inverse
for $y\ge1$ by
\begin{equation}\label{eq:N-def}
 N(y)=\min\{n\ge2:D_n\ge y\}.
\end{equation}
This is a staircase, not a smooth function.  Its asymptotic relation to the
inverse gamma core is best encoded by the locations of its jumps.

\subsection{Core-coordinate thresholds}

Define
\begin{equation}\label{eq:alpha-n-def}
 \alpha_n=\InvCore(D_n),
 \qquad\text{equivalently}\qquad
 \Core(\alpha_n)=D_n.
\end{equation}
Because $\InvCore$ is increasing,
\begin{equation}\label{eq:N-threshold}
 \boxed{\quad
 N(y)=\min\{n\ge2:\InvCore(y)\le\alpha_n\}.
 \quad}
\end{equation}
The thresholds $\alpha_n$ are exponentially close to the integers.

\begin{theorem}[Threshold displacement]\label{thm:threshold-displacement}
Let $s_n=(-1)^n$, $J_n=\Tail(n)$, and $P_n=\psi(n+1)$.  Then
\begin{align}
 \alpha_n-n={}&
 \frac{s_nJ_n}{\Core'(n)}
 -\frac{\Core''(n)J_n^2}{2\Core'(n)^3}\notag\\
 &+\frac{s_n\bigl(3\Core''(n)^2-\Core'(n)\Core'''(n)\bigr)J_n^3}
 {6\Core'(n)^5}
 +\Ord\!\left(\frac{J_n^4}{\Core(n)^4P_n}\right).
 \label{eq:alpha-n-Taylor}
\end{align}
In particular,
\begin{equation}\label{eq:alpha-n-leading}
 \boxed{\quad
 \alpha_n=n+(-1)^n
 \frac{\e\Tail(n)}{n!\psi(n+1)}
 +\Ord\!\left(\frac1{(n!)^2n^2\log n}\right).
 \quad}
\end{equation}
Using the Bell expansion of $\Tail(n)$ yields a complete factorially small
threshold transseries.
\end{theorem}

\begin{proof}
From \eqref{eq:integer-correction},
$D_n=\Core(n)+s_nJ_n$.  Taylor-expand the local inverse of $\Core$ at
$\Core(n)$.  The standard inverse-function derivatives are
\begin{align*}
 (\Core^{-1})'&=\frac1{\Core'},\\
 (\Core^{-1})''&=-\frac{\Core''}{\Core'^3},\\
 (\Core^{-1})'''&=\frac{3\Core''^2-\Core'\Core'''}{\Core'^5}.
\end{align*}
Substitution of the increment $s_nJ_n$ gives
\eqref{eq:alpha-n-Taylor}.  Since
$\Core'(n)=(n!/\e)P_n$, $J_n=\Ord(n^{-1})$, and
$\Core''(n)=\Core(n)(P_n^2+P_n')$, the first term is
\eqref{eq:alpha-n-leading} and the second is of the stated factorial order.
\end{proof}

The first term can be written entirely in elementary integer data because
$P_n=H_n-\gamma$ and
\begin{equation}\label{eq:J-n-exact}
 \Tail(n)=(-1)^n\left(D_n-\frac{n!}{\e}\right).
\end{equation}
For asymptotic work, however, the Bell expansion is more revealing:
\begin{equation}\label{eq:alpha-n-Bell}
 \alpha_n-n\sim
 \frac{(-1)^n\e}{n!(H_n-\gamma)}
 \left(\frac1n-\frac2{n^2}+\frac5{n^3}-\frac{15}{n^4}+\cdots\right)
 +\text{higher factorial sectors}.
\end{equation}

\subsection{Why the staircase has no ordinary smooth transseries}

The difference $N(y)-\InvCore(y)$ remains bounded, but it does not approach
a limit as $y$ moves continuously through successive intervals
$(D_{n-1},D_n]$.  In fact it traverses essentially a full unit interval
between jumps.  Hence there is no expansion
\begin{equation}\label{eq:no-smooth-staircase}
 N(y)=\InvCore(y)+c_0+\littleo(1)
\end{equation}
uniformly for all real $y\to\infty$ with a fixed constant $c_0$.  The
correct decomposition is
\begin{equation}\label{eq:staircase-decomp}
 N(y)=\left\lceil\InvCore(y)-\vartheta(y)\right\rceil,
\end{equation}
where the threshold correction $\vartheta$ is supported through the tiny
shifts $\alpha_n-n$.  Formula \eqref{eq:N-threshold}, rather than a smooth
Poincar\'e series, is the exact asymptotic normal form.

\begin{remark}[Away from thresholds]
If $\InvCore(y)$ stays a distance much larger than
$1/(n!n\log n)$ from every integer $n$, then the exponentially small shifts
in \eqref{eq:alpha-n-leading} do not affect the result and
$N(y)=\lceil\InvCore(y)\rceil$.  Only in exponentially narrow windows near
integer core coordinates does the parity-dependent derangement correction
change the step location.
\end{remark}

\section{Endpoint-Mellin tails and coefficient sequences}
\label{sec:endpoint-Mellin}

The Bell numbers arise from one particular endpoint kernel, but the
mechanism is much broader.  Let $\varphi$ be analytic in a neighborhood of
$t=1$ and define
\begin{equation}\label{eq:Jphi}
 J_\varphi(x)=\int_0^1t^x\varphi(t)\dd t.
\end{equation}
With $t=\e^{-s}$,
\begin{equation}\label{eq:Jphi-Laplace}
 J_\varphi(x)=\int_0^\infty\e^{-(x+1)s}
 g(s)\dd s,
 \qquad g(s)=\varphi(\e^{-s}).
\end{equation}

\begin{theorem}[Endpoint-jet expansion]\label{thm:endpoint-jet}
Suppose $g$ is analytic at $s=0$ and has sufficient growth control for
Watson's lemma.  Then
\begin{equation}\label{eq:Jphi-Watson}
 J_\varphi(x)\sim\sum_{r=0}^{\infty}
 \frac{g^{(r)}(0)}{(x+1)^{r+1}}.
\end{equation}
Equivalently, after re-expansion in powers of $x^{-1}$,
\begin{equation}\label{eq:Jphi-x-expansion}
 J_\varphi(x)\sim\sum_{k=1}^{\infty}\frac{\mu_k}{x^k},
 \qquad
 \mu_k=\sum_{r=0}^{k-1}(-1)^{k-r-1}\binom{k-1}{r}g^{(r)}(0).
\end{equation}
\end{theorem}

\begin{proof}
Insert the Taylor series
$g(s)=\sum_{r\ge0}g^{(r)}(0)s^r/r!$ into
\eqref{eq:Jphi-Laplace} and integrate termwise asymptotically:
\[
 \int_0^\infty\e^{-(x+1)s}s^r\dd s=\frac{r!}{(x+1)^{r+1}}.
\]
This proves \eqref{eq:Jphi-Watson}.  Expanding
$(x+1)^{-r-1}=x^{-r-1}(1+x^{-1})^{-r-1}$ and collecting the coefficient of
$x^{-k}$ gives \eqref{eq:Jphi-x-expansion}.
\end{proof}

For the subfactorial,
\begin{equation}\label{eq:sub-phi-g}
 \varphi(t)=\e^{t-1},
 \qquad
 g(s)=\e^{\e^{-s}-1},
 \qquad
 g^{(r)}(0)=(-1)^r\Bell_r.
\end{equation}
Thus \eqref{eq:Jphi-Watson} becomes the shifted Bell expansion
\eqref{eq:J-shifted-Bell}.  Other endpoint kernels generate other
coefficient sequences.  For example:

\begin{itemize}
\item $\varphi(t)=1$ gives $J_\varphi(x)=1/(x+1)$ and a geometric sequence;
\item $\varphi(t)=t^\beta$ gives $1/(x+\beta+1)$;
\item $\varphi(t)=\exp(\alpha t)$ gives coefficients governed by Touchard
      polynomials $T_r(\alpha)$;
\item products or compositions in $\varphi$ produce complete Bell
      polynomials in its endpoint derivatives.
\end{itemize}

\subsection{Inverse transseries for gamma-type endpoint perturbations}

Consider
\begin{equation}\label{eq:gamma-endpoint-family}
 F(x)=\Core(x)+\sum_{q=1}^{Q}
 \e^{\ii\omega_q x}J_{\varphi_q}(x),
\end{equation}
where $\Core$ is a rapid core and the frequencies are fixed.  Combining
\cref{thm:operator-sectors,thm:endpoint-jet} yields the nested expansion
\begin{equation}\label{eq:gamma-endpoint-inverse}
 F^{-1}(y)\sim a(y)+\sum_{n\ge1}\frac1{y^n}
 \sum_{\omega\in n\Omega}
 \e^{\ii\omega a(y)}C_{n,\omega}(a(y)),
\end{equation}
where each $C_{n,\omega}$ has an inverse-power expansion determined by the
endpoint jets of the $\varphi_q$ and by the logarithmic derivative of the
core.  The subfactorial is the one-frequency case
$Q=1$, $\omega_1=\sigma\pi$, and
$\varphi_1(t)=\e^{t-1}$.

\section{Further structure of the rapid-core calculus}
\label{sec:further-structure}

\subsection{Functoriality under exact changes of core coordinate}

The derivation of \cref{thm:master-reversion} shows that the entire calculus
is functorial under the coordinate change $u=\Core(x)$.  Algebraically, the
inverse problem factors as
\begin{equation}\label{eq:factorization}
 (\Core+\Res)^{-1}
 =a\circ(\operatorname{id}+\Res\circ a)^{-1}.
\end{equation}
This factorization separates two tasks:

\begin{enumerate}
\item reverse the near-identity map $u\mapsto u+\rho(u)$;
\item transport the result back through the exact core inverse $a$.
\end{enumerate}

The coefficient formula \eqref{eq:master-LB} is precisely the transport of
ordinary near-identity reversion through this factorization.  It therefore
inherits associativity and composition laws from the usual group of formal
diffeomorphisms tangent to the identity.

\subsection{Bell polynomials inside Bell sectors}

There are two unrelated but interacting appearances of ``Bell'' in this
problem.

\begin{enumerate}
\item The numerical Bell numbers $\Bell_k$ arise from the endpoint kernel
      and constitute the coefficients of $\Tail(x)$.
\item Fa\`a di Bruno and repeated differentiation in
      \eqref{eq:Cn-sigma} organize products of derivatives through partial
      Bell polynomials.
\end{enumerate}

To make the second appearance explicit, write
\begin{equation}\label{eq:rho-y}
 \rho(y)=\Res(a(y)).
\end{equation}
Then
\begin{equation}\label{eq:LB-Bell-polynomial}
 \frac{\dd^{n-1}}{\dd y^{n-1}}\bigl(a'(y)\rho(y)^n\bigr)
 =\sum_{j=0}^{n-1}\binom{n-1}{j}a^{(j+1)}(y)
 \frac{\dd^{n-1-j}}{\dd y^{n-1-j}}\rho(y)^n.
\end{equation}
The last derivative can be expressed through exponential partial Bell
polynomials in $\rho',\rho'',\ldots$.  Thus the endpoint Bell numbers enter
as data, while combinatorial Bell polynomials enter as the universal chain
rule of inversion.

\subsection{Multiple algebraic tails}

Suppose
\begin{equation}\label{eq:R-multiscale}
 \Res(x)\sim\sum_{q\ge q_0}r_q(\log x)x^{-q}
\end{equation}
with polynomial or rational logarithmic coefficients.  At outer sector $n$,
$\Res(x)^n$ begins at $x^{-nq_0}$.  The operator
$\Lam^{-1}D_x-r$ preserves the leading algebraic order when
$\Lam$ grows logarithmically or more rapidly.  Consequently,
\begin{equation}\label{eq:sector-leading-order-q0}
 X_n(x)=\Ord\!\left(
 \frac{x^{-nq_0}}{\Lam(x)}
 \right).
\end{equation}
This gives the support of the inverse transseries before any coefficient is
computed.

\subsection{Cores of exponential form}

If
\begin{equation}\label{eq:A-exp-Phi}
 \Core(x)=\exp(\Phi(x)),
\end{equation}
then $\Lam=\Phi'$, and the master coefficients simplify to
\begin{equation}\label{eq:Xn-Phi}
 X_n(x)=\frac{(-1)^n}{n!}
 \left(\frac{D_x}{\Phi'}-(n-1)\right)\cdots
 \left(\frac{D_x}{\Phi'}-1\right)
 \left(\frac{\Res(x)^n}{\Phi'(x)}\right).
\end{equation}
Gamma-type functions fit this form with
$\Phi(x)=\log\Gamma(x+1)-1$.  Barnes-type and generalized hyperfactorial
cores also fit it after taking logarithms; only the preliminary inversion of
$\Phi$ changes.  The perturbative sector calculus remains identical.

\subsection{A universality principle}

\begin{proposition}[Universal leading sector law]
\label{prop:universal-sector}
Let $(\Core,\Res)$ be admissible and suppose, for every fixed $n$,
\[
 D_x\log(\Res^n/\Lam)=\littleo(\Lam).
\]
Then
\begin{equation}\label{eq:universal-sector-law}
 F^{-1}(y)-a(y)
 \sim-\frac1{\Lam(a(y))}
 \sum_{n\ge1}\frac1n
 \left(\frac{\Res(a(y))}{y}\right)^n
\end{equation}
in the sense of matching the leading coefficient in each fixed outer
sector.  Equivalently, the sector leaders are generated by the logarithm in
\eqref{eq:frozen-log}.
\end{proposition}

\begin{proof}
Under the stated hypothesis, every $D_x/\Lam$ contribution in
\eqref{eq:Xn-general} is lower order.  The remaining composition is
$(-1)^{n-1}(n-1)!$, proving \eqref{eq:Xn-leading-general}.  Summation over
$n$ is the Taylor series of $\log(1-z)$.
\end{proof}

\subsection{When the logarithmic resummation fails}

The hypothesis of \cref{prop:universal-sector} is essential.  It fails if
the tail contains phases whose frequencies grow as fast as $\Lam$, boundary
layers whose derivatives amplify the coefficient, or rapidly varying
singularities near the inverse point.  In such cases $D_x/\Lam$ and
multiplication by $-r$ contribute at the same order, and the leading
all-sector sum is no longer logarithmic.  The exact operator formula
\eqref{eq:Xn-general} remains valid and should be used instead.

\section{Numerical validation}
\label{sec:numerics}

The numerical tests in this section use high-precision arithmetic.  The core
inverse $a=\InvCore(y)$ is obtained by Newton iteration applied to
\begin{equation}\label{eq:numeric-core-equation}
 \log\Gamma(a+1)-1-\log y=0,
\end{equation}
starting from the Lambert seed $u_0-1/2$.  The principal lateral inverse is
then obtained from Newton iteration on
\begin{equation}\label{eq:numeric-sub-equation}
 \frac{\Gamma(z+1)}\e+\e^{\ii\pi z}\Tail(z)-y=0,
\end{equation}
starting from the first-sector approximation.  The endpoint integral is
evaluated directly; the partial-fraction series \eqref{eq:J-meromorphic}
provides an independent check.

\subsection{Accuracy of the centered inverse-gamma expansion}

Let $a^{[M]}$ be the truncation containing $p_1,\ldots,p_M$, with
$a^{[0]}=u_0-1/2$.  \Cref{tab:gamma-core-errors} reports
$|a^{[M]}-a|$.

\begin{table}[htbp]
\centering
\small
\caption{Absolute errors in the centered inverse-gamma expansion.}
\label{tab:gamma-core-errors}
\resizebox{\textwidth}{!}{%
\begin{tabular}{@{}r l rrrrr@{}}
\toprule
$y$ & exact $a=\InvCore(y)$
& $M=0$ & $M=1$ & $M=2$ & $M=3$ & $M=4$\\
\midrule
$10^3$
& $6.69048431391006103$
& $2.9344\times 10^{-3}$ & $4.7978\times 10^{-6}$ & $2.7535\times 10^{-8}$
& $3.6454\times 10^{-10}$ & $9.2245\times 10^{-12}$\\
$10^6$
& $9.87684651274886097$
& $1.7154\times 10^{-3}$ & $1.2715\times 10^{-6}$ & $3.5241\times 10^{-9}$
& $2.2938\times 10^{-11}$ & $2.8457\times 10^{-13}$\\
$10^{10}$
& $13.5617402632380633$
& $1.1206\times 10^{-3}$ & $4.3640\times 10^{-7}$ & $6.6159\times 10^{-10}$
& $2.3767\times 10^{-12}$ & $1.6235\times 10^{-14}$\\
$10^{50}$
& $41.5613050474933070$
& $2.6492\times 10^{-4}$ & $1.0626\times 10^{-8}$ & $1.8221\times 10^{-12}$
& $7.5138\times 10^{-16}$ & $5.8501\times 10^{-19}$\\
\bottomrule
\end{tabular}%
}
\end{table}

The consistent reduction by several orders of magnitude confirms both the
half-shifted normal form and the rational functions
\eqref{eq:p1}--\eqref{eq:p4}.  For still larger $M$, the Stirling series is
asymptotic rather than convergent; the optimal truncation and any desired
hyperasymptotic continuation should be handled through the exact residual
\eqref{eq:gamma-residual}.

\subsection{Accuracy of the Bell sectors}

Let
\begin{align}
 X_+^{[0]}(y)&=a(y),\label{eq:X0-num}\\
 X_+^{[1]}(y)&=a(y)+\frac{\e^{\ii\pi a(y)}}yC_{1,+}(a(y)),
 \label{eq:X1-num}\\
 X_+^{[2]}(y)&=X_+^{[1]}(y)
 +\frac{\e^{2\ii\pi a(y)}}{y^2}C_{2,+}(a(y)).
 \label{eq:X2-num}
\end{align}
\Cref{tab:lateral-errors} compares these approximations with a direct
high-precision root of \eqref{eq:numeric-sub-equation}.

\begin{table}[htbp]
\centering
\small
\caption{Principal lateral inverse: exact displacement and sector errors.}
\label{tab:lateral-errors}
\begin{tabular}{@{}r >{\raggedright\arraybackslash}p{5.4cm} rrr@{}}
\toprule
$y$ & $X_+(y)-a(y)$
& $|X_+^{[0]}-X_+|$
& $|X_+^{[1]}-X_+|$
& $|X_+^{[2]}-X_+|$\\
\midrule
$10^3$
& $3.32599835278648\times 10^{-5}-4.87639893172798\times 10^{-5}\ii$
& $5.9027\times 10^{-5}$ & $1.1643\times 10^{-8}$ & $3.4746\times 10^{-12}$\\
$10^6$
& $-3.35445523390485\times 10^{-8}+1.36671789756972\times 10^{-8}\ii$
& $3.6222\times 10^{-8}$ & $4.4488\times 10^{-15}$ & $8.2808\times 10^{-22}$\\
$10^{10}$
& $-4.70356185958560\times 10^{-13}+2.39449487665544\times 10^{-12}\ii$
& $2.4403\times 10^{-12}$ & $2.0481\times 10^{-23}$ & $2.6085\times 10^{-34}$\\
\bottomrule
\end{tabular}
\end{table}

The ratios are consistent with one extra factor of approximately
$1/(ya)$ for each added sector.  At $y=10^{10}$, two Bell sectors improve
the core approximation from about twelve decimal orders in the displacement
to roughly thirty-four decimal orders in absolute error.

\subsection{Residual-based verification}

A coefficient formula can be checked without knowing the exact inverse.
Substitute a truncation $G_N(y)$ into the defining equation and expand
\begin{equation}\label{eq:formal-residual-check}
 \Sub_\sigma(G_N(y))-y.
\end{equation}
The coefficients through outer order $y^{-N+1}$ must vanish, and the first
nonzero residual must be of order $y^{-N}a^{-N-1}$ up to slow logarithmic
factors.  Division by
$\Sub_\sigma'(G_N)\sim y\log a$ then gives the inverse error predicted in
\eqref{eq:sub-all-orders}.  This residual test was used independently of the
operator derivation for the first two sectors.

\section{Implementation details and symbolic workflow}
\label{sec:implementation}

\subsection{Stable numerical evaluation}

For large $y$, direct evaluation of $\Gamma(a+1)/\e-y$ can overflow even
when the residual is small.  The core equation should be solved in logarithmic
form, as in \eqref{eq:numeric-core-equation}.  A robust procedure is:

\begin{algorithm}[Numerical principal inverse subfactorial]
\label{alg:numeric-sub}
Given $y>0$ sufficiently large and a working precision of $p$ digits:
\begin{enumerate}
\item Compute $\lambda=\log(\e y/\sqrt{2\pi})$,
      $w=W_0(\lambda/\e)$, and $a_0=\lambda/w-1/2$.
\item Optionally add $p_1/u_0$, $p_2/u_0^3$, and further centered
      corrections.
\item Apply Newton iteration to
\[
 h(a)=\log\Gamma(a+1)-1-\log y,
 \qquad h'(a)=\psi(a+1),
\]
      to obtain the exact core inverse $a$ at the desired precision.
\item Evaluate $\Tail(a)$ either by the integral
      $\int_0^1\e^{t-1}t^a\dd t$ or by
      \eqref{eq:J-meromorphic}.  For large positive $a$, a short optimally
      truncated Bell expansion with its signed remainder bound is also
      effective.
\item Form the first-sector guess
\[
 z_0=a-\frac{\e^{\ii\pi a}\Tail(a)}{y\psi(a+1)}.
\]
\item Apply complex Newton iteration to
\[
 H(z)=\Sub_+(z)-y,
\]
      using
\[
 H'(z)=\frac{\Gamma(z+1)}\e\psi(z+1)
 +\e^{\ii\pi z}\bigl(\Tail'(z)+\ii\pi\Tail(z)\bigr).
\]
\end{enumerate}
Because the starting error is already of second Bell-sector order, very few
Newton steps are normally required.
\end{algorithm}

\subsection{Symbolic construction of the slow coefficients}

The two recurrences in this article are triangular in different senses.

\begin{description}[leftmargin=2.6em,style=nextline]
\item[Core recurrence.]
Equation \eqref{eq:pm-recurrence} determines $p_m(w)$ from
$p_1,\ldots,p_{m-1}$ by coefficient extraction in the dummy variable
$r=u_0^{-1}$.

\item[Outer-sector recurrence.]
Equation \eqref{eq:Cn-sigma} determines the complete $n$th Bell sector by
$n-1$ applications of a first-order differential operator.  It does not
require solving a new implicit equation.
\end{description}

A computer algebra implementation should represent the slow coefficient
field as a truncated differential algebra generated by
\begin{equation}\label{eq:slow-generators}
 x^{-1},\quad (\log x)^{-1},\quad P(x),\quad P'(x),\ldots,
 \quad \Tail(x),\quad\Tail'(x),\ldots,
\end{equation}
and only afterward substitute the Bell and Bernoulli expansions.  This
prevents expression swell and preserves exact identities such as
\eqref{eq:C2-exact}.

\subsection{Pseudocode}

\begin{lstlisting}
# Core coefficients p_m(w)
Delta = 0
for m in 1..M:
    E = w*Delta + (1/r + Delta)*log(1 + r*Delta)
    E += sum(c[j]*r^(2*j-1)*(1+r*Delta)^(-(2*j-1)), j=1..m)
    p[m] = -coeff(E, r, 2*m-1)/(1+w)
    Delta += p[m]*r^(2*m-1)

# Bell sectors for the lateral inverse
for n in 1..N:
    H = J(x)^n / P(x)
    for r in 1..n-1:
        H = diff(H,x)/P(x) + (I*sigma*n*pi/P(x) - r)*H
    C[n] = (-1)^n * H / factorial(n)

X = a(y) + sum(exp(I*sigma*n*pi*a(y))*C[n](a(y))/y^n, n=1..N)
\end{lstlisting}

\subsection{Precision management}

The correction $X_+(y)-a(y)$ is of order $1/(ya\log a)$ and may be far
smaller than the magnitude of $a$.  If that correction is required to $p$
decimal digits relative to itself, the core $a$ must be computed with
roughly
\begin{equation}\label{eq:guard-digits}
 p+\log_{10}y+\log_{10}a+\log_{10}\log a
\end{equation}
absolute guard digits beyond the decimal point.  Alternatively, solve
directly for the displacement $\delta=z-a$ using the normalized equation
\begin{equation}\label{eq:normalized-delta}
 \frac{\Core(a+\delta)}{\Core(a)}-1
 +\frac{\e^{\ii\pi(a+\delta)}\Tail(a+\delta)}{y}=0,
\end{equation}
which avoids subtracting two nearly equal large numbers.

\section{Interpretation of the transseries hierarchy}
\label{sec:hierarchy}

\subsection{Outer and inner asymptotics}

The exact sector formula \eqref{eq:sub-all-orders} is an asymptotic expansion
in the outer parameter $1/y$ with coefficients depending on the slow
variable $a(y)$.  Each coefficient has its own expansion in $1/a$ and
$1/\log a$.  Symbolically,
\begin{equation}\label{eq:two-level-symbol}
 X_\sigma(y)-a(y)
 \sim\sum_{n\ge1}y^{-n}\e^{\ii\sigma n\pi a(y)}
 \sum_{p\ge n}\sum_{q\ge1}c_{n,p,q}^{(\sigma)}
 a(y)^{-p}(\log a(y))^{-q}.
\end{equation}
The lower bound $p\ge n$ follows because $\Tail^n$ begins with $a^{-n}$.
At fixed $n$ and $p$, only finitely many logarithmic structures occur at a
given truncation order.

\subsection{Ordering relative to the core expansion}

The algebraic Stirling correction $u_0^{-(2m-1)}$ is much larger than any
fixed Bell sector $y^{-n}$ as $y\to\infty$.  Hence every finite number of
Stirling terms precedes the first Bell sector.  The centered Stirling series
is divergent at sufficiently high order, so a fully hyperasymptotic ordering
can contain exponentially improved gamma-core remainders.  The exact-core
form \eqref{eq:sub-all-orders} deliberately avoids asserting an artificial
global order between those hyperasymptotic core contributions and the Bell
sectors.  It says instead:

\begin{quote}
Resolve the gamma core to whatever accuracy is required, then add the
subfactorial-specific sectors as an independent perturbative hierarchy.
\end{quote}

This modular viewpoint is particularly important at astronomically large
arguments, where the exponentially improved remainder of a deeply truncated
Stirling expansion may cross the scale $y^{-1}$.

\subsection{Phase as a chirp in logarithmic time}

With $Y=\log y$,
\begin{equation}\label{eq:a-Y-leading}
 a(y)\sim\frac{Y}{W_0(Y/\e)}.
\end{equation}
The principal phase is therefore
\begin{equation}\label{eq:chirp-phase}
 \e^{\ii\pi a(y)}
 \sim\exp\left(
 \ii\pi\frac{Y}{W_0(Y/\e)}
 \right).
\end{equation}
Its instantaneous frequency with respect to $Y$ is
\begin{equation}\label{eq:chirp-frequency}
 \frac{\dd}{\dd Y}\bigl(\pi a(\e^Y)\bigr)
 =\frac{\pi}{P(a(\e^Y))}
 \sim\frac{\pi}{\log a(\e^Y)}.
\end{equation}
Thus the phase oscillates indefinitely but ever more slowly in logarithmic
time $Y$.  With respect to $y$ itself, its derivative is smaller by another
factor $1/y$.

\subsection{Stokes and branch data}

The lateral sign $\sigma$ is not an arbitrary decoration: it records how the
negative argument in the incomplete gamma function is approached.  It is
therefore part of the analytic data of the problem.  The two transseries are
complex conjugates on the positive real $y$ axis, and their median gives a
real interpolation only after the forward functions are averaged.  Choosing
one continuation and silently taking the real part of its inverse would not
produce the inverse of any canonically specified forward function.

\section{Extensions, conjectures, and research directions}
\label{sec:extensions}

\subsection{General winding branches}

For the branch \eqref{eq:winding-branch} with
$\omega_k=(2k+1)\pi$, all formulas remain valid after replacing
$\sigma\pi$ by $\omega_k$.  In particular,
\begin{equation}\label{eq:Cn-k}
 C_{n,k}(x)=\frac{(-1)^n}{n!}
 \prod_{r=n-1}^{1}
 \left[
 \frac{D_x}{P}+\frac{\ii n\omega_k}{P}-r
 \right]
 \left(\frac{\Tail(x)^n}{P(x)}\right),
\end{equation}
where the product is composition from right to left.  At fixed $k$, the
same error estimate holds.  Uniformity as $|k|\to\infty$ is a separate
multiple-scale problem because the phase derivative
$|\omega_k|/P(x)$ need no longer be small.

\subsection{Transseries for parameter-dependent derangements}

Many generalized derangement numbers have exponential generating functions
of the form
\begin{equation}\label{eq:generalized-derangement-egf}
 \frac{\exp(-\alpha z)}{(1-z)^\beta}
\end{equation}
or finite products of such factors.  Their dominant continuation is often a
gamma-type core multiplied by powers, while the truncation error is an
endpoint-Mellin integral.  The two-stage method developed here applies:
first invert the dominant logarithmic phase, then use
\eqref{eq:Xn-general} for the endpoint correction.  Touchard, Bell,
Stirling, or generalized Bell polynomials will appear according to the
endpoint kernel.

\subsection{Large-order behavior of the centered coefficients}

The coefficients $p_m(w)$ inherit the factorial divergence of the Stirling
series.  The initial examples suggest the structured form
\begin{equation}\label{eq:pm-structure}
 p_m(w)=\frac{(-1)^{m+1}P_m(w)}{d_m(1+w)^{2m-1}},
 \qquad \deg P_m=2m-2,
\end{equation}
with $P_m$ having positive coefficients for the computed orders.

\begin{conjecture}[Positivity pattern for centered inverse-gamma coefficients]
\label{conj:pm-positive}
After clearing the minimal positive integer denominator, the polynomial
$(-1)^{m+1}(1+w)^{2m-1}p_m(w)$ has nonnegative coefficients for every
$m\ge1$ and is strictly positive for $w\ge0$.
\end{conjecture}

A proof would likely require a sign-preserving reformulation of the
triangular recurrence \eqref{eq:pm-recurrence}; the alternating Bernoulli
input makes coefficientwise positivity far from automatic.

\subsection{Large-order behavior of Bell sectors}

At fixed $x$, the Lagrange series in the bookkeeping parameter has a radius
determined by the nearest critical point of
$\Core(x)+\varepsilon\Res(x)$.  As $n\to\infty$ with $y$ fixed, the
coefficients $C_{n,\sigma}$ will therefore show a different large-order law
from the fixed-$n$, $y\to\infty$ asymptotics developed here.  Locating the
nearest coalescing inverse branches would determine that law and the optimal
outer-sector truncation.

\subsection{Resurgent completion}

The centered Stirling series for the gamma core and the outer Bell-sector
series have different sources of divergence.  A full resurgent analysis
would combine:

\begin{enumerate}
\item Borel singularities of the log-gamma Stirling series;
\item singularities in the Lagrange parameter caused by critical points of
      the perturbed inverse problem;
\item the incomplete-gamma branch data encoded by $\sigma$;
\item possible scale crossings between exponentially improved core terms and
      $y^{-n}$ Bell sectors.
\end{enumerate}

The exact-core normal form \eqref{eq:sub-all-orders} is a convenient starting
point because it cleanly separates the first item from the remaining three.

\subsection{Formalization prospects}

The core pieces are well suited to formal proof assistants:

\begin{itemize}
\item the exact series \eqref{eq:J-meromorphic} and remainder identity
      \eqref{eq:J-Bell-rem} reduce to dominated convergence, a finite
      geometric identity, and Dobi\'nski's formula;
\item the master inversion formula \eqref{eq:master-LB} can be developed in
      a filtered differential algebra of formal series;
\item the operator recurrence \eqref{eq:Xn-general} is finite at every
      requested sector order;
\item residual certification \eqref{eq:residual-to-error} separates formal
      coefficient cancellation from analytic error bounds.
\end{itemize}

For the ProveIt corpus, a natural first milestone is a formal theorem for the
Bell remainder bound and a purely algebraic implementation of the
valuation-raising operators.  The analytic existence theorem can then be
added using established gamma and polygamma estimates.

\section{Conclusion}
\label{sec:conclusion}

The asymptotic inverse of the subfactorial is not obtained by merely replacing
$n!$ with a continuous factorial and applying one Lambert $W$.  The complete
structure is
\begin{equation}\label{eq:conclusion-structure}
 \boxed{
 \text{inverse gamma core}
 \quad+\quad
 \text{Lambert--Stirling corrections}
 \quad+\quad
 \text{oscillatory Bell sectors in }y^{-1}.
 }
\end{equation}
The exact decomposition
\[
 \Sub_\sigma(z)=\Gamma(z+1)/\e+
 \e^{\ii\sigma\pi z}\Tail(z)
\]
turns the problem into a rapid-core perturbation.  The endpoint tail
$\Tail$ has Bell coefficients with an explicit signed remainder, and
Lagrange--B\"urmann inversion transports those coefficients into a locally
finite hierarchy of inverse sectors.  The distinguished principal branch is
complex, its lower lateral companion is its conjugate, and a real inverse
requires a separately specified real interpolation.  The discrete inverse is
a staircase whose asymptotic information is concentrated in factorially
small shifts of its jump thresholds.

The reusable result is the operator calculus
\[
 X_n(x)=\frac{(-1)^n}{n!}
 \left(\frac{D_x}{\Lam}-(n-1)\right)\cdots
 \left(\frac{D_x}{\Lam}-1\right)
 \left(\frac{\Res(x)^n}{\Lam(x)}\right),
\]
with phase-adapted variants for oscillatory tails.  It applies whenever a
function consists of a rapidly varying invertible core plus a tail small in
relative size.  The subfactorial supplies a particularly clean model because
its tail is both exactly representable and combinatorially structured.

\appendix

\section{Initial Bell data and endpoint coefficients}
\label{app:Bell-data}

The first Bell numbers are
\begin{center}
\begin{tabular}{@{}c|rrrrrrrrrrrr@{}}
\toprule
$k$ & 0&1&2&3&4&5&6&7&8&9&10&11\\
\midrule
$\Bell_k$
&1&1&2&5&15&52&203&877&4140&21147&115975&678570\\
\bottomrule
\end{tabular}
\end{center}
Hence
\begin{align}
 \Tail(x)\sim{}&x^{-1}-2x^{-2}+5x^{-3}-15x^{-4}+52x^{-5}
 -203x^{-6}\notag\\
 &+877x^{-7}-4140x^{-8}+21147x^{-9}-115975x^{-10}
 +678570x^{-11}-\cdots.
 \label{eq:J-long}
\end{align}
The shifted form is
\begin{align}
 \Tail(x)\sim{}&(x+1)^{-1}-(x+1)^{-2}+2(x+1)^{-3}
 -5(x+1)^{-4}\notag\\
 &+15(x+1)^{-5}-52(x+1)^{-6}+203(x+1)^{-7}-\cdots.
 \label{eq:J-long-shifted}
\end{align}

\section{Derivation of the second inverse sector}
\label{app:C2-derivation}

From \eqref{eq:Cn-sigma},
\[
 C_{2,\sigma}=\frac12\mathcal L^{(\sigma)}_{1,2}
 \left(\frac{\Tail^2}{P}\right).
\]
By definition,
\begin{align*}
 \mathcal L^{(\sigma)}_{1,2}
 \left(\frac{\Tail^2}{P}\right)
 &=\frac1P\left(\frac{\Tail^2}{P}\right)'
 +\left(\frac{2\ii\sigma\pi}{P}-1\right)
 \frac{\Tail^2}{P}\\
 &=\frac{2\Tail\Tail'}{P^2}
 -\frac{\Tail^2P'}{P^3}
 +\frac{2\ii\sigma\pi\Tail^2}{P^2}
 -\frac{\Tail^2}{P}.
\end{align*}
Division by two gives \eqref{eq:C2-exact}.  This calculation also shows how
phase differentiation and slow differentiation enter at different powers of
$1/P$.

\section{First terms of the third sector}
\label{app:C3}

The exact compact formula is
\begin{equation}\label{eq:C3-operator}
 C_{3,\sigma}(x)=-\frac16
 \mathcal L^{(\sigma)}_{2,3}
 \mathcal L^{(\sigma)}_{1,3}
 \left(\frac{\Tail(x)^3}{P(x)}\right).
\end{equation}
Without expanding the full rational expression, its leading slow terms are
\begin{equation}\label{eq:C3-leading}
 C_{3,\sigma}(x)
 =-\frac1{3x^3\log x}
 +\Ord\!\left(\frac1{x^3(\log x)^2}
              +\frac1{x^4\log x}\right).
\end{equation}
Thus the first three outer sectors begin
\begin{align}
 X_\sigma(y)-a(y)={}&
 -\frac{\e^{\ii\sigma\pi a}}{ya\log a}
 -\frac{\e^{2\ii\sigma\pi a}}{2y^2a^2\log a}
 -\frac{\e^{3\ii\sigma\pi a}}{3y^3a^3\log a}
 +\text{lower slow orders}.
 \label{eq:first-three-leaders}
\end{align}
These are precisely the first three terms in the logarithmic resummation
\eqref{eq:sub-frozen-log} after replacing $\Tail(a)$ by $1/a$ and $P(a)$ by
$\log a$.

\section{A direct formal proof of local finiteness}
\label{app:local-finiteness}

Let $\mathscr K$ be a differential coefficient field in the variable $x$
containing $\Lam^{\pm1}$ and $\Res$, and consider formal sums
\begin{equation}\label{eq:formal-ring}
 \mathscr K((y^{-1}))_{\ge0}
 =\left\{\sum_{n\ge0}y^{-n}f_n(a(y)):f_n\in\mathscr K\right\}.
\end{equation}
Define the outer valuation of a nonzero series by
\begin{equation}\label{eq:outer-valuation}
 \nu\left(\sum_{n\ge n_0}y^{-n}f_n(a(y))\right)=n_0.
\end{equation}
Equation \eqref{eq:valuation-derivative} implies
\begin{equation}\label{eq:valuation-gain}
 \nu(D_yS)\ge\nu(S)+1.
\end{equation}
The $n$th Lagrange term in \eqref{eq:master-LB} starts with one explicit
factor $y^{-1}$ and contains $n-1$ derivatives.  It therefore has valuation
at least $n$.  Consequently, the coefficient of $y^{-N}$ receives
contributions only from Lagrange terms with $n\le N$.  This is local
finiteness, and it justifies all formal coefficient extractions without any
analytic convergence assumption.

\section{Notation crosswalk}
\label{app:notation}

\begin{longtable}{@{}p{0.18\textwidth}p{0.72\textwidth}@{}}
\toprule
Symbol & Meaning\\
\midrule
\endfirsthead
\toprule
Symbol & Meaning\\
\midrule
\endhead
$D_n={!n}$ & integer derangement number or subfactorial\\
$\Sub_\sigma(z)$ & lateral incomplete-gamma continuation, $\arg(-1)=\sigma\pi$\\
$\Sub_{\med}(z)$ & real median interpolation $(\Sub_++\Sub_-)/2$\\
$\Core(x)$ & gamma core $\Gamma(x+1)/\e$ in the subfactorial application; a general rapid core elsewhere\\
$\Tail(x)$ & endpoint tail $\e^{-1}\int_0^1\e^tt^x\dd t$\\
$\InvCore(y)$ & large real inverse of $\Core$, equal to $\Gamma_0^{-1}(\e y)-1$\\
$P(x)$ & $\psi(x+1)$, the logarithmic derivative of the gamma core\\
$\Lam(x)$ & general logarithmic derivative $\Core'(x)/\Core(x)$\\
$X_\sigma(y)$ & distinguished large inverse branch of $\Sub_\sigma$\\
$C_{n,\sigma}(x)$ & slow coefficient of the $n$th lateral Bell sector\\
$\lambda$ & $\log(\e y/\sqrt{2\pi})$\\
$w$ & $W_0(\lambda/\e)$\\
$u_0$ & Lambert seed $\lambda/w$\\
$p_m(w)$ & centered inverse-gamma Stirling coefficient\\
$N(y)$ & discrete generalized inverse $\min\{n:D_n\ge y\}$\\
$\alpha_n$ & core-coordinate threshold $\InvCore(D_n)$\\
\bottomrule
\end{longtable}

\begin{thebibliography}{99}

\bibitem{BorweinCorless2018}
J.~M. Borwein and R.~M. Corless,
\newblock Gamma and factorial in the Monthly,
\newblock \emph{Amer. Math. Monthly} \textbf{125} (2018), no.~5, 400--424,
\newblock \url{https://doi.org/10.1080/00029890.2018.1420983}.

\bibitem{Comtet1974}
L.~Comtet,
\newblock \emph{Advanced Combinatorics},
\newblock Reidel, Dordrecht, 1974.

\bibitem{CorlessAmenyouJeffrey2017}
R.~M. Corless, F.~K. Amenyou, and D.~J. Jeffrey,
\newblock Properties and computation of the functional inverse of gamma,
\newblock in \emph{2017 19th International Symposium on Symbolic and Numeric
Algorithms for Scientific Computing (SYNASC)}, 2017, pp.~65--72,
\newblock \url{https://doi.org/10.1109/SYNASC.2017.00020}.

\bibitem{CorlessEtAl1996}
R.~M. Corless, G.~H. Gonnet, D.~E.~G. Hare, D.~J. Jeffrey, and D.~E. Knuth,
\newblock On the Lambert $W$ function,
\newblock \emph{Adv. Comput. Math.} \textbf{5} (1996), 329--359,
\newblock \url{https://doi.org/10.1007/BF02124750}.

\bibitem{DLMFGamma}
NIST Digital Library of Mathematical Functions,
\newblock Chapter~5: Gamma function, especially \S5.11,
\newblock \url{https://dlmf.nist.gov/5.11}, accessed 3 September 2026.

\bibitem{DLMFIncompleteGamma}
NIST Digital Library of Mathematical Functions,
\newblock Chapter~8: Incomplete gamma and related functions,
\newblock \url{https://dlmf.nist.gov/8}, accessed 3 September 2026.

\bibitem{FlajoletSedgewick2009}
P.~Flajolet and R.~Sedgewick,
\newblock \emph{Analytic Combinatorics},
\newblock Cambridge University Press, Cambridge, 2009.

\bibitem{Hassani2020}
M.~Hassani,
\newblock Derangements and alternating sum of permutations by integration,
\newblock \emph{J. Integer Seq.} \textbf{23} (2020), Article 20.7.8.

\bibitem{Olver1997}
F.~W.~J. Olver,
\newblock \emph{Asymptotics and Special Functions},
\newblock AKP Classics, A K Peters, Wellesley, MA, 1997.

\bibitem{ParisKaminski2001}
R.~B. Paris and D.~Kaminski,
\newblock \emph{Asymptotics and Mellin--Barnes Integrals},
\newblock Cambridge University Press, Cambridge, 2001.

\bibitem{Pedersen2015}
H.~L. Pedersen,
\newblock Inverses of gamma functions,
\newblock \emph{Constr. Approx.} \textbf{41} (2015), no.~2, 251--267,
\newblock \url{https://doi.org/10.1007/s00365-014-9239-1}.

\bibitem{ProveItSeries}
V.~Reshetnikov,
\newblock ProveIt: series and transseries draft collection,
\newblock \url{https://github.com/VladimirReshetnikov/ProveIt/tree/main/Analysis/FabiusFunction/docs/semi-formalized-research-frontiers/drafts/series-and-transseries},
\newblock accessed 3 September 2026.

\bibitem{Uchiyama2012}
M.~Uchiyama,
\newblock The principal inverse of the gamma function,
\newblock \emph{Proc. Amer. Math. Soc.} \textbf{140} (2012), no.~4,
1343--1348,
\newblock \url{https://doi.org/10.1090/S0002-9939-2011-11023-2}.

\bibitem{vanDerHoeven2006}
J.~van der Hoeven,
\newblock \emph{Transseries and Real Differential Algebra},
\newblock Lecture Notes in Mathematics 1888, Springer, Berlin, 2006.

\end{thebibliography}

\end{document}
